% =====================================================================
%  Chapter 10: Double Integrals
%  Sections 10.1 and 10.2
% =====================================================================

\chapter{Double Integrals}
\label{ch:10}

Chapter \ref{ch:9} asked pointwise questions.  We looked for peaks, valleys, saddle
points, and constrained best points.  The tools were local:
\[
    df,\qquad \nabla f,\qquad H_f.
\]
They told us what happens near one point.

Now the question changes.  A farmer does not only ask where a field is tallest.  She
asks how much crop the field will produce.  A builder does not only ask where a roof is
highest.  He asks how much air lies under it.  A metal plate is not only hot at one
point.  It stores heat across its whole surface.

Single-variable integration added tiny pieces along an interval.  A double integral adds
tiny pieces across a region in the plane.

\section{From area to volume by slicing}
\label{sec:10.1}

The single-variable integral began with area under a curve.  A double integral begins
with the same idea, but the base is no longer an interval.  The base is a region in the
plane, and the graph above it is a surface.

\subsection{Area as a limit of sums}

Imagine a garden bed whose edge is curved.  The owner wants its area, but the shape is
not a rectangle, a triangle, or a circle.  One practical method is to lay a square grid over
the garden and count squares.

Large squares give a rough estimate.  Smaller squares give a better one.

\begin{figure}[h]
\centering
\begin{tikzpicture}[scale=0.72, >=stealth]
    \draw[step=1, gray!25, very thin] (0,0) grid (8,5);

    \draw[thick, fill=green!10]
        plot[smooth cycle] coordinates {
            (1.0,0.8)
            (2.1,0.25)
            (4.3,0.55)
            (6.8,1.2)
            (7.25,2.7)
            (6.0,4.25)
            (3.6,4.65)
            (1.35,3.75)
            (0.55,2.1)
        };

    \node at (4,-0.7) {a curved region estimated by grid squares};
\end{tikzpicture}
\caption{Area can be estimated by adding the areas of many small squares.}
\end{figure}

If each small square has area
\[
    \Delta A,
\]
and \(N\) of them lie mostly inside the garden, then
\[
    \operatorname{Area} \approx N\Delta A.
\]
This is only an estimate, because boundary squares are awkward.  But as the grid gets
finer, the boundary error gets squeezed into a thinner and thinner strip.

The idea is not new.  In single-variable calculus, we approximated area under a curve by
rectangles:
\[
    \sum f(x_i^*)\,\Delta x.
\]
Now each small piece has area
\[
    \Delta A
\]
instead of length \(\Delta x\).

For a plain area problem, every small piece contributes just its own area.  So area is
the limit of sums like
\[
    \sum \Delta A.
\]
Equivalently,
\[
\boxed{
    \operatorname{Area}(R)=\iint_R 1\,dA.
}
\]
This formula says: add the constant value \(1\) over every tiny area piece in \(R\).

\subsection{Volume under a surface}

Now put a surface above the region.  Suppose a slanted canopy covers a rectangular
market stall:
\[
    0\le x\le 4,\qquad 0\le y\le 3.
\]
Its height in meters is
\[
    f(x,y)=2+\frac{x}{4}+\frac{y}{6}.
\]
We want the volume of air under the canopy.

Cut the floor into \(1\)-meter by \(1\)-meter squares.  Over each small square, pretend
the canopy height is constant and equal to the height at the center of that square.
Then each piece is approximately a rectangular column:
\[
    \text{volume of one column}
    \approx
    \text{height at center}\times \text{base area}.
\]

The centers have
\[
    x=0.5,1.5,2.5,3.5
\]
and
\[
    y=0.5,1.5,2.5.
\]
Since each small square has area \(1\), the volume estimate is
\[
    \sum f(x_i^*,y_j^*)\cdot 1.
\]
Compute the sum:
\[
\begin{aligned}
\sum f(x_i^*,y_j^*)
&=
\sum\left(2+\frac{x_i^*}{4}+\frac{y_j^*}{6}\right)\\
&=
12\cdot 2
+
\frac14\left(3(0.5+1.5+2.5+3.5)\right)\\
&\qquad
+
\frac16\left(4(0.5+1.5+2.5)\right)\\
&=
24+\frac14(24)+\frac16(18)\\
&=
24+6+3\\
&=
33.
\end{aligned}
\]
So the estimated volume is
\[
    33\text{ m}^3.
\]

\begin{figure}[h]
\centering
\begin{tikzpicture}
    \begin{axis}[
        view={45}{35},
        width=13cm, height=9cm,
        axis lines=center,
        xlabel={\(x\)}, ylabel={\(y\)}, zlabel={\(z\)},
        xmin=0, xmax=4.8, ymin=0, ymax=3.8, zmin=0, zmax=4,
        xtick={0,1,2,3,4}, ytick={0,1,2,3}, ztick={0,1,2,3,4},
        ticklabel style={font=\footnotesize},
        enlargelimits=0.05
    ]
        % Corrected Painter's Algorithm: Draw from back-to-front 
        % Since the camera is at +x, -y, we draw large y first and small x first.
        \foreach \y in {2,1,0} {
            \foreach \x in {0,1,2,3} {
                \pgfmathsetmacro{\h}{2 + (\x+0.5)/4 + (\y+0.5)/6}
                \pgfmathsetmacro{\xp}{\x+1}
                \pgfmathsetmacro{\yp}{\y+1}
                
                \edef\temp{
                    % Front Face (facing the viewer, constant y)
                    \noexpand\filldraw[fill=cyan!40!white, draw=cyan!80!black, thin]
                        (\x,\y,0) -- (\xp,\y,0) -- (\xp,\y,\h) -- (\x,\y,\h) -- cycle;
                    
                    % Right Face (facing the viewer, constant x)
                    \noexpand\filldraw[fill=cyan!60!white, draw=cyan!80!black, thin]
                        (\xp,\y,0) -- (\xp,\yp,0) -- (\xp,\yp,\h) -- (\xp,\y,\h) -- cycle;
                    
                    % Top Face (facing up, constant z)
                    \noexpand\filldraw[fill=cyan!10!white, draw=cyan!80!black, thin]
                        (\x,\y,\h) -- (\xp,\y,\h) -- (\xp,\yp,\h) -- (\x,\yp,\h) -- cycle;
                }
                \temp 
            }
        }

        % Plot the exact sample points at the center of each column top
        \addplot3[
            only marks,
            mark=*,
            mark size=1.0pt,
            color=black
        ] coordinates {
            (0.5,0.5,2.208) (1.5,0.5,2.458) (2.5,0.5,2.708) (3.5,0.5,2.958)
            (0.5,1.5,2.375) (1.5,1.5,2.625) (2.5,1.5,2.875) (3.5,1.5,3.125)
            (0.5,2.5,2.542) (1.5,2.5,2.792) (2.5,2.5,3.042) (3.5,2.5,3.292)
        };

        % Translucent canopy surface
        \addplot3[
            surf,
            domain=0:4,
            domain y=0:3,
            samples=2,
            fill=red!40,
            draw=red!60!black,
            opacity=0.4,
            thick
        ] {2 + x/4 + y/6};
        
    \end{axis}
\end{tikzpicture}
\caption{A double integral adds many small column volumes.}
\end{figure}

The same pattern works for any surface
\[
    z=f(x,y)
\]
above a region \(R\).  Divide \(R\) into many small pieces.  On each piece, choose a
sample point \((x^*,y^*)\).  The small column has approximate volume
\[
    f(x^*,y^*)\,\Delta A.
\]
Add all the columns:
\[
    \sum f(x^*,y^*)\,\Delta A.
\]
If these sums settle down to one number as the pieces shrink, that number is the
volume under the surface.

\subsection{Riemann sums over rectangles}

The market canopy computation had the same shape as a one-variable Riemann sum,
but with area pieces instead of length pieces.

For a rectangle
\[
    R=[a,b]\times[c,d],
\]
we cut the \(x\)-interval into pieces and the \(y\)-interval into pieces.  This creates
many small subrectangles.

A typical subrectangle has side lengths
\[
    \Delta x_i,\qquad \Delta y_j,
\]
so its area is
\[
    \Delta A_{ij}=\Delta x_i\,\Delta y_j.
\]
Choose a sample point
\[
    (x_{ij}^*,y_{ij}^*)
\]
inside that subrectangle.  Then the Riemann sum is
\[
\boxed{
    \sum_i\sum_j f(x_{ij}^*,y_{ij}^*)\,\Delta A_{ij}.
}
\]

This expression has a simple meaning:

\[
\begin{array}{c|c}
\text{factor} & \text{meaning} \\ \hline
f(x_{ij}^*,y_{ij}^*) & \text{height, density, temperature, or value on the piece}\\[1mm]
\Delta A_{ij} & \text{area of the piece}\\[1mm]
f(x_{ij}^*,y_{ij}^*)\Delta A_{ij} & \text{small contribution}
\end{array}
\]

When the function is a height, the contribution is a small volume.  When the function
is a density, the contribution is a small mass.  When the function is temperature, the
contribution is temperature times area, which leads to average temperature after we
divide by total area.

\begin{definition}[Double integral, first meaning]
If the Riemann sums
\[
    \sum_i\sum_j f(x_{ij}^*,y_{ij}^*)\,\Delta A_{ij}
\]
approach one number as the subrectangles get smaller, independent of the sample
points, then that number is called the \emph{double integral} of \(f\) over \(R\), written
\[
\boxed{
    \iint_R f(x,y)\,dA.
}
\]
\end{definition}

The symbol \(dA\) means a tiny positive area piece.

There is also a differential-forms way to write the same object.  In Chapter
\ref{ch:2}, the form
\[
    dx\wedge dy
\]
measured signed area in the \(xy\)-plane.  On the usual orientation of the plane, a tiny
rectangle whose sides point in the positive \(x\)- and \(y\)-directions has signed area
\[
    dx\wedge dy=\Delta x\,\Delta y.
\]
So the oriented version of the same integral is
\[
\boxed{
    \iint_R f(x,y)\,dx\wedge dy.
}
\]
For ordinary area, volume, mass, and average value, we usually write
\[
    f\,dA.
\]
When orientation matters, especially near Green's theorem, the wedge notation will
return.

\subsection{What a double integral measures}

A double integral is not tied to one picture.  It is an adding machine for quantities
spread over a region.

\[
\begin{array}{c|c|c}
\text{Function} & \text{Units of } f & \text{Meaning of } \displaystyle\iint_R f\,dA \\ \hline
f(x,y)=1 & 1 & \text{area of }R\\[1mm]
f(x,y)=\text{height} & \text{meters} & \text{volume under a surface}\\[1mm]
\rho(x,y)=\text{density} & \text{kg/m}^2 & \text{mass of a plate}\\[1mm]
T(x,y)=\text{temperature} & ^\circ\text{C} & \text{temperature-area total}\\[1mm]
q(x,y)=\text{charge density} & \text{C/m}^2 & \text{total charge}
\end{array}
\]

\begin{example}[Mass of a thin plate]
A thin metal plate occupies a region \(R\) in the plane.  Its density at \((x,y)\) is
\[
    \rho(x,y)\quad \text{kg/m}^2.
\]
A tiny area piece \(\Delta A\) near \((x,y)\) has mass approximately
\[
    \rho(x,y)\Delta A.
\]
Adding all such pieces gives
\[
\boxed{
    M=\iint_R \rho(x,y)\,dA.
}
\]
The double integral is not a volume here.  It is mass.
\end{example}

\begin{example}[Signed volume can cancel]
Let
\[
    f(x,y)=x-y
\]
on the unit square
\[
    R=[0,1]\times[0,1].
\]
Above the diagonal \(x=y\), the function is positive.  Below the diagonal, it is
negative.

For every point \((x,y)\), the reflected point \((y,x)\) has value
\[
    f(y,x)=y-x=-(x-y).
\]
The square is symmetric under this reflection.  So the positive and negative
contributions cancel:
\[
\boxed{
    \iint_R (x-y)\,dA=0.
}
\]
This does not mean the surface is flat.  It means the signed volume above the square
equals the signed volume below it.
\end{example}

The idea is now visible: choose small area pieces, multiply each by the value sitting on
it, and add.  But the phrase ``small area pieces'' still needs a careful meaning.  If two
people choose different sample points, will they get the same answer?  If the function
jumps around wildly, can the sums refuse to settle?  The next section turns the picture
of columns into a definition strong enough to compute with.

\section{Double integrals over rectangles}
\label{sec:10.2}

Section \ref{sec:10.1} built the double integral from a picture: columns over small
rectangles.  Now we make the pieces precise.  The goal is the same as in
single-variable calculus: define the integral in a way that does not depend on a lucky
choice of sample points.

\subsection{Partitions and sample points}

Start with a rectangle
\[
    R=[a,b]\times[c,d].
\]
A partition of the \(x\)-interval is a list
\[
    a=x_0<x_1<\cdots<x_m=b.
\]
A partition of the \(y\)-interval is a list
\[
    c=y_0<y_1<\cdots<y_n=d.
\]
Together they cut \(R\) into subrectangles
\[
    R_{ij}=[x_{i-1},x_i]\times[y_{j-1},y_j].
\]
The area of \(R_{ij}\) is
\[
    \Delta A_{ij}=\Delta x_i\,\Delta y_j,
\]
where
\[
    \Delta x_i=x_i-x_{i-1},
    \qquad
    \Delta y_j=y_j-y_{j-1}.
\]

\begin{figure}[h]
\centering
\begin{tikzpicture}[scale=0.75, >=stealth]
    \draw[->] (-0.4,0) -- (6.6,0) node[right] {\(x\)};
    \draw[->] (0,-0.4) -- (0,4.7) node[above] {\(y\)};

    \draw[thick] (0,0) rectangle (6,4);

    \foreach \x in {1,2.4,4.1} {
        \draw[gray!50] (\x,0) -- (\x,4);
    }
    \foreach \y in {0.9,2.2,3.1} {
        \draw[gray!50] (0,\y) -- (6,\y);
    }

    \foreach \p in {
        (0.55,0.45),(1.65,0.55),(3.2,0.35),(5.0,0.65),
        (0.4,1.5),(1.7,1.35),(3.3,1.75),(5.2,1.45),
        (0.75,2.75),(1.5,2.6),(3.5,2.55),(5.0,2.75),
        (0.55,3.55),(1.8,3.55),(3.3,3.65),(5.2,3.5)
    } {
        \fill \p circle (1.5pt);
    }

    \node at (3,-0.8) {one sample point in each subrectangle};
\end{tikzpicture}
\caption{A partition cuts a rectangle into smaller rectangles.  A Riemann sum chooses
one sample point in each piece.}
\end{figure}

Choose a sample point
\[
    (x_{ij}^*,y_{ij}^*)\in R_{ij}
\]
for each subrectangle.  The corresponding Riemann sum is
\[
\boxed{
    S=
    \sum_{i=1}^{m}\sum_{j=1}^{n}
    f(x_{ij}^*,y_{ij}^*)\,\Delta x_i\,\Delta y_j.
}
\]

\begin{example}[A midpoint estimate on a rectangle]
Let
\[
    R=[0,4]\times[0,2],
\]
and let
\[
    f(x,y)=3+x+\frac12y.
\]
Cut \(R\) into four rectangles of size \(2\times1\).  The midpoint sample points are
\[
    (1,0.5),\quad (3,0.5),\quad (1,1.5),\quad (3,1.5).
\]
Each subrectangle has area
\[
    \Delta A=2.
\]
The midpoint Riemann sum is
\[
\begin{aligned}
S
&=
2\Bigl[
f(1,0.5)+f(3,0.5)+f(1,1.5)+f(3,1.5)
\Bigr]\\
&=
2\Bigl[
\left(3+1+\frac14\right)
+\left(3+3+\frac14\right)
+\left(3+1+\frac34\right)
+\left(3+3+\frac34\right)
\Bigr]\\
&=
2\left[
\frac{17}{4}+\frac{25}{4}+\frac{19}{4}+\frac{27}{4}
\right]\\
&=
2\cdot 22\\
&=
44.
\end{aligned}
\]
So this grid estimates
\[
    \iint_R f\,dA
\]
by \(44\).
\end{example}

For this linear function, the midpoint estimate is already exact.  We will prove such
facts later using iterated integrals.  For now, the sum is the object to understand.

\subsection{Integrable functions, informally}

A function is integrable over a rectangle when the Riemann sums settle down to one
number as the rectangles get small.

The sample points should stop mattering.  The exact shape of the grid should stop
mattering.  Only the function and the region should matter.

\begin{definition}[Integrable over a rectangle]
A bounded function
\[
    f:R\to\mathbb R
\]
on a rectangle
\[
    R=[a,b]\times[c,d]
\]
is \emph{integrable over \(R\)} if all Riemann sums
\[
    \sum_i\sum_j f(x_{ij}^*,y_{ij}^*)\,\Delta A_{ij}
\]
approach the same number as the largest side length of the subrectangles approaches
\(0\).

That common number is
\[
\boxed{
    \iint_R f\,dA.
}
\]
\end{definition}

There is another way to say the same thing.  On each subrectangle, the function has a
largest possible value and a smallest possible value, at least for the functions we first
meet.  If we use all the largest values, we get an upper estimate.  If we use all the
smallest values, we get a lower estimate.

If the upper and lower estimates can be forced close together by using a fine enough
partition, then the function is integrable.

\begin{theorem}[Continuous functions on rectangles are integrable]
If
\[
    f
\]
is continuous on a rectangle
\[
    R=[a,b]\times[c,d],
\]
then
\[
    f
\]
is integrable over \(R\).
\end{theorem}

\begin{proof}[Idea of proof]
A continuous function on a closed rectangle cannot oscillate too much inside very
small rectangles.  If the partition is fine enough, then the largest and smallest values
of \(f\) on each subrectangle are close.  Multiplying those small differences by the
areas of the subrectangles and adding them gives upper and lower sums that are close.
So every reasonable Riemann sum is trapped between two nearly equal numbers.
\end{proof}

We required continuity because it gives a clean guarantee.  Without it, the conclusion
can fail.

\begin{example}[A bounded function that is not integrable]
Define a function on the unit square
\[
    R=[0,1]\times[0,1]
\]
by
\[
    D(x,y)=
    \begin{cases}
    1, & x\text{ and }y\text{ are both rational},\\
    0, & \text{otherwise}.
    \end{cases}
\]
Every small rectangle contains points with rational coordinates and also points where
at least one coordinate is irrational.  So on every subrectangle,
\[
    \text{largest value}=1,
    \qquad
    \text{smallest value}=0.
\]
Every upper sum is therefore
\[
    1\cdot \operatorname{Area}(R)=1,
\]
and every lower sum is
\[
    0.
\]
The upper and lower sums never get close.

The function is bounded, but not integrable.  It jumps too wildly on every tiny piece.
\end{example}

The boundedness condition is also part of this definition.

\begin{example}[An unbounded function needs a different treatment]
Define
\[
    u(x,y)=
    \begin{cases}
    \dfrac{1}{x^2+y^2}, & (x,y)\ne(0,0),\\[2mm]
    0, & (x,y)=(0,0),
    \end{cases}
\]
on
\[
    [-1,1]\times[-1,1].
\]
The function is unbounded near the origin.  A sample point very close to the origin
can make a Riemann sum extremely large.

This function does not fit the ordinary rectangle definition above.  Integrals of
unbounded functions require a limiting process.  Those are improper double integrals,
which return in Section \ref{sec:10.7}.
\end{example}

The rectangle itself also has finite area.  If we try to integrate
\[
    f(x,y)=1
\]
over the whole plane, the sums do not approach a finite number.  They grow with the
area being covered.  Infinite regions also need improper integrals.

For now, rectangles are finite, functions are bounded, and continuity is the safe
condition that makes the integral exist.

\subsection{Average value over a rectangle}

Single-variable calculus defined the average value of a function on an interval by
\[
    f_{\text{avg}}
    =
    \frac{1}{b-a}\int_a^b f(x)\,dx.
\]
The denominator was the length of the interval.

For a rectangle, length is replaced by area.

\begin{definition}[Average value over a rectangle]
Let \(f\) be integrable over a rectangle \(R\) with positive area.  The \emph{average
value} of \(f\) over \(R\) is
\[
\boxed{
    f_{\text{avg}}
    =
    \frac{1}{\operatorname{Area}(R)}
    \iint_R f\,dA.
}
\]
\end{definition}

The units work out exactly as they should.  If \(f\) is temperature in degrees Celsius,
then
\[
    \iint_R f\,dA
\]
has units
\[
    ^\circ\text{C}\cdot \text{m}^2.
\]
Dividing by area leaves
\[
    ^\circ\text{C}.
\]

\begin{example}[Average height of the market canopy]
In Section \ref{sec:10.1}, the canopy over the rectangle
\[
    [0,4]\times[0,3]
\]
had estimated volume
\[
    33\text{ m}^3.
\]
The floor area is
\[
    4\cdot 3=12\text{ m}^2.
\]
So the average canopy height is
\[
    \frac{33}{12}=\frac{11}{4}=2.75\text{ m}.
\]
The average height is the height of a flat roof over the same floor that would enclose
the same volume.
\end{example}

That last sentence is often the best way to remember average value:

\[
\boxed{
    \text{average height}\times\text{base area}
    =
    \text{total volume}.
}
\]

\subsection{Numerical estimates from tables and contour maps}

In applications, we often do not know a formula for \(f(x,y)\).  We may only have
measurements.

Suppose a greenhouse floor is
\[
    0\le x\le 6,\qquad 0\le y\le 4,
\]
and temperature is measured at the centers of six equal rectangles.  Each rectangle has
dimensions
\[
    2\times2,
\]
so each has area
\[
    4.
\]

\[
\begin{array}{c|ccc}
 & x=1 & x=3 & x=5 \\ \hline
y=3 & 18 & 19 & 20\\
y=1 & 22 & 24 & 23
\end{array}
\]

The midpoint estimate for the integral is
\[
\begin{aligned}
\iint_R T\,dA
&\approx
4(18+19+20+22+24+23)\\
&=
4(126)\\
&=
504.
\end{aligned}
\]
The total area is
\[
    6\cdot4=24.
\]
So the estimated average temperature is
\[
    T_{\text{avg}}
    \approx
    \frac{504}{24}
    =
    21^\circ\text{C}.
\]

This is a double integral computed from a table.  The table supplies the sample values.
The grid supplies the area pieces.

Contour maps give another kind of estimate.  They do not give values at isolated
points; they show curves of equal value.  We can estimate by splitting the region into
bands.

\begin{figure}[h]
\centering
\begin{tikzpicture}[scale=0.82, >=stealth]
    \draw[thick, fill=blue!6]
        plot[smooth cycle] coordinates {
            (0.4,1.2)
            (1.2,0.3)
            (3.2,0.2)
            (5.3,0.9)
            (5.7,2.4)
            (4.7,3.8)
            (2.5,4.2)
            (0.7,3.3)
        };

    \draw[blue!55!black, thick]
        plot[smooth cycle] coordinates {
            (1.0,1.4)
            (1.7,0.9)
            (3.4,0.8)
            (4.7,1.3)
            (4.9,2.4)
            (4.0,3.2)
            (2.4,3.4)
            (1.2,2.7)
        };
    \node[blue!55!black] at (5.35,1.35) {\(1\) m};

    \draw[blue!65!black, thick]
        plot[smooth cycle] coordinates {
            (1.6,1.7)
            (2.2,1.25)
            (3.4,1.25)
            (4.1,1.8)
            (4.1,2.5)
            (3.4,2.9)
            (2.3,2.9)
            (1.7,2.4)
        };
    \node[blue!65!black] at (4.55,2.55) {\(2\) m};

    \draw[blue!75!black, thick]
        plot[smooth cycle] coordinates {
            (2.25,1.9)
            (2.8,1.65)
            (3.35,1.8)
            (3.55,2.35)
            (3.1,2.55)
            (2.45,2.45)
        };
    \node[blue!75!black] at (3.75,2.05) {\(3\) m};

    \node at (3,-0.55) {depth contours for a small pond};
\end{tikzpicture}
\caption{A contour map can be turned into an estimate by using the areas between
contours.}
\end{figure}

\begin{example}[Estimating water volume from depth bands]
A pond is divided by contour lines into depth bands.  Suppose the estimated areas are:

\[
\begin{array}{c|c|c}
\text{depth band} & \text{typical depth} & \text{area} \\ \hline
0\text{ to }1\text{ m} & 0.5\text{ m} & 120\text{ m}^2\\
1\text{ to }2\text{ m} & 1.5\text{ m} & 200\text{ m}^2\\
2\text{ to }3\text{ m} & 2.5\text{ m} & 160\text{ m}^2\\
3\text{ to }4\text{ m} & 3.5\text{ m} & 60\text{ m}^2
\end{array}
\]

Then the water volume is estimated by
\[
\begin{aligned}
V
&\approx
0.5(120)+1.5(200)+2.5(160)+3.5(60)\\
&=
60+300+400+210\\
&=
970\text{ m}^3.
\end{aligned}
\]
The pond area is
\[
    120+200+160+60=540\text{ m}^2.
\]
So the estimated average depth is
\[
    \frac{970}{540}\approx 1.80\text{ m}.
\]
\end{example}

Both the table method and the contour method are Riemann sums in disguise.  Each
small contribution has the form
\[
    \text{value}\times\text{area}.
\]

So far, our double integrals have been built by estimates.  That is honest, but slow.
Single-variable calculus became powerful when area sums turned into ordinary
antiderivatives.  The same thing happens here in a surprising way: a double integral
over a rectangle can often be computed by doing two single-variable integrals, one
after the other.  The next step is to slice the rectangle into strips and let the old
Fundamental Theorem of Calculus do part of the work.
% =====================================================================
%  Chapter 10: Double Integrals
%  Sections 10.3 and 10.4
% =====================================================================

\section{Iterated integrals and Fubini's theorem}
\label{sec:10.3}

Section \ref{sec:10.2} defined a double integral over a rectangle as a limit of sums:
\[
    \iint_R f(x,y)\,dA.
\]
That definition tells us what the integral means.  It adds tiny contributions
\[
    f(x,y)\,\Delta A
\]
over many small rectangles.

But the definition is not always the best way to compute.  In single-variable calculus,
we did not evaluate every definite integral by taking a limit of Riemann sums.  Once the
integral was defined, the Fundamental Theorem of Calculus gave us a practical method.

Something similar happens here.  A double integral can often be computed by doing two
ordinary one-variable integrals, one after the other.

\subsection{A roof sliced two ways}

A market stall has a rectangular base
\[
    0\le x\le 4,
    \qquad
    0\le y\le 3.
\]
The height of its slanted roof above the floor is modeled by
\[
    h(x,y)=5+x+2y.
\]
The volume under the roof is
\[
    V=\iint_R h(x,y)\,dA,
\]
where
\[
    R=[0,4]\times[0,3].
\]

Instead of adding tiny columns all at once, freeze \(x\).  Then \(y\) runs from \(0\) to
\(3\), and the roof above that vertical strip has cross-sectional area
\[
    A(x)=\int_0^3 (5+x+2y)\,dy.
\]
Compute the inner integral:
\[
\begin{aligned}
    A(x)
    &=
    \int_0^3 (5+x+2y)\,dy      \\
    &=
    \Bigl[(5+x)y+y^2\Bigr]_0^3 \\
    &=
    3(5+x)+9                  \\
    &=
    24+3x.
\end{aligned}
\]
So a slice at position \(x\) has area
\[
    24+3x.
\]
Now add those slice areas from \(x=0\) to \(x=4\):
\[
\begin{aligned}
    V
    &=
    \int_0^4 A(x)\,dx       \\
    &=
    \int_0^4 (24+3x)\,dx    \\
    &=
    \Bigl[24x+\frac32x^2\Bigr]_0^4 \\
    &=
    96+24                  \\
    &=
    120.
\end{aligned}
\]
Thus
\[
\boxed{
    \iint_R (5+x+2y)\,dA=120.
}
\]

The same volume can be sliced the other way.  Freeze \(y\), let \(x\) run from \(0\) to
\(4\), and compute
\[
    B(y)=\int_0^4 (5+x+2y)\,dx.
\]
Then
\[
\begin{aligned}
    B(y)
    &=
    \Bigl[5x+\frac12x^2+2yx\Bigr]_0^4 \\
    &=
    20+8+8y                         \\
    &=
    28+8y.
\end{aligned}
\]
Now add these horizontal slice areas:
\[
\begin{aligned}
    V
    &=
    \int_0^3 B(y)\,dy              \\
    &=
    \int_0^3 (28+8y)\,dy           \\
    &=
    \Bigl[28y+4y^2\Bigr]_0^3       \\
    &=
    84+36                          \\
    &=
    120.
\end{aligned}
\]

Both slicing methods give the same volume:
\[
\boxed{
    \int_0^4\int_0^3 (5+x+2y)\,dy\,dx
    =
    \int_0^3\int_0^4 (5+x+2y)\,dx\,dy.
}
\]

\begin{figure}[h]
\centering
\begin{tikzpicture}[scale=0.9, >=stealth]
    \draw[->] (-0.3,0) -- (5.1,0) node[right] {\(x\)};
    \draw[->] (0,-0.3) -- (0,4.1) node[above] {\(y\)};

    \draw[fill=gray!8] (0,0) rectangle (4,3);

    \draw[fill=blue!20, draw=blue!60!black, thick] (1.35,0) rectangle (1.75,3);
    \node[blue!60!black] at (2.75,2.4) {fix \(x\)};
    \draw[->, blue!60!black, thick] (2.45,2.25) -- (1.75,2.25);

    \draw[fill=red!18, draw=red!65!black, thick, opacity=0.75] (0,1.05) rectangle (4,1.45);
    \node[red!65!black] at (3.0,0.65) {fix \(y\)};
    \draw[->, red!65!black, thick] (2.55,0.75) -- (2.55,1.05);

    \node[below] at (4,0) {\(4\)};
    \node[left] at (0,3) {\(3\)};
    \node[below left] at (0,0) {\(0\)};
    \node at (2,-0.7) {\(R=[0,4]\times[0,3]\)};
\end{tikzpicture}
\caption{A rectangle can be added by vertical slices or by horizontal slices.}
\end{figure}

The picture is simple, but the idea is important.  A double integral over a rectangle can
be computed by reducing it to ordinary integrals.

\subsection{Iterated integrals}

The expression
\[
    \int_0^4\int_0^3 (5+x+2y)\,dy\,dx
\]
means this:

First, hold \(x\) fixed and integrate with respect to \(y\):
\[
    \int_0^3 (5+x+2y)\,dy.
\]
Then integrate the result with respect to \(x\):
\[
    \int_0^4 \left(\int_0^3 (5+x+2y)\,dy\right) dx.
\]

This is called an \emph{iterated integral}, because one integral is nested inside another.

\begin{definition}[Iterated integral over a rectangle]
Let
\[
    R=[a,b]\times[c,d].
\]
The iterated integral
\[
\boxed{
    \int_a^b\int_c^d f(x,y)\,dy\,dx
}
\]
means
\[
\boxed{
    \int_a^b
    \left(
        \int_c^d f(x,y)\,dy
    \right)
    dx.
}
\]
In the inner integral, \(x\) is held constant and \(y\) is the variable.

Similarly,
\[
\boxed{
    \int_c^d\int_a^b f(x,y)\,dx\,dy
}
\]
means
\[
\boxed{
    \int_c^d
    \left(
        \int_a^b f(x,y)\,dx
    \right)
    dy.
}
\]
\end{definition}

The inner integral makes a slice.  The outer integral adds the slices.

A small notation warning is useful here.  The symbol
\[
    dy\,dx
\]
in an iterated integral is an instruction for the order of calculation:
\[
    \text{integrate in }y\text{ first, then in }x.
\]
It is not a wedge product.

With the standard orientation of the \(xy\)-plane, the corresponding area \(2\)-form is
\[
    dx\wedge dy.
\]
So, for a positively oriented rectangle,
\[
    \iint_R f\,dA
    =
    \int_R f\,dx\wedge dy.
\]
But
\[
    dy\wedge dx=-dx\wedge dy.
\]
Thus the wedge product remembers orientation, while the ordinary area element
\[
    dA
\]
remembers positive area.

\subsection{Fubini's theorem for rectangles}

The roof example showed that the two orders gave the same answer.  Fubini's theorem
tells us when this is guaranteed.

\begin{theorem}[Fubini's theorem for rectangles]
Let
\[
    R=[a,b]\times[c,d],
\]
and suppose \(f\) is continuous on \(R\).  Then
\[
\boxed{
    \iint_R f(x,y)\,dA
    =
    \int_a^b\int_c^d f(x,y)\,dy\,dx
    =
    \int_c^d\int_a^b f(x,y)\,dx\,dy.
}
\]
\end{theorem}

The continuity assumption is a safe hypothesis.  It guarantees that the double integral
exists and that slicing in either order gives the same total.

\begin{proof}[Idea of the proof]
For each fixed \(x\), define
\[
    A(x)=\int_c^d f(x,y)\,dy.
\]
Geometrically, \(A(x)\) is the signed cross-sectional area of the surface over the vertical
line segment above \(x\).  Adding those cross-sections gives
\[
    \int_a^b A(x)\,dx
    =
    \int_a^b\int_c^d f(x,y)\,dy\,dx.
\]
This adds the same tiny rectangular contributions that the double integral adds:
\[
    f(x,y)\,\Delta x\,\Delta y.
\]

Similarly, if
\[
    B(y)=\int_a^b f(x,y)\,dx,
\]
then
\[
    \int_c^d B(y)\,dy
    =
    \int_c^d\int_a^b f(x,y)\,dx\,dy
\]
adds the same tiny contributions, only grouped into horizontal slices instead of vertical
ones.

Both methods refine to the same Riemann sum limit, so both give the double integral.
\end{proof}

\subsection{A separated example}

Sometimes the integrand splits into an \(x\)-part and a \(y\)-part.  Then the double
integral factors.

Let
\[
    R=[0,3]\times[0,\pi],
\]
and let
\[
    f(x,y)=(1+x^2)(2+\sin y).
\]
Then
\[
\begin{aligned}
    \iint_R (1+x^2)(2+\sin y)\,dA
    &=
    \int_0^3\int_0^\pi (1+x^2)(2+\sin y)\,dy\,dx.
\end{aligned}
\]
Since \(1+x^2\) is constant with respect to \(y\),
\[
\begin{aligned}
    \int_0^\pi (1+x^2)(2+\sin y)\,dy
    &=
    (1+x^2)\int_0^\pi (2+\sin y)\,dy    \\
    &=
    (1+x^2)\Bigl[2y-\cos y\Bigr]_0^\pi \\
    &=
    (1+x^2)(2\pi+2).
\end{aligned}
\]
Now integrate in \(x\):
\[
\begin{aligned}
    \iint_R (1+x^2)(2+\sin y)\,dA
    &=
    \int_0^3 (1+x^2)(2\pi+2)\,dx\\
    &=
    (2\pi+2)\Bigl[x+\frac{x^3}{3}\Bigr]_0^3\\
    &=
    (2\pi+2)(3+9)\\
    &=
    24\pi+24.
\end{aligned}
\]
Equivalently,
\[
\boxed{
    \iint_R (1+x^2)(2+\sin y)\,dA
    =
    \left(\int_0^3 (1+x^2)\,dx\right)
    \left(\int_0^\pi (2+\sin y)\,dy\right).
}
\]

The double integral factors because the rectangle has independent limits and the
integrand is a product:
\[
    f(x,y)=g(x)h(y).
\]

\subsection{Choosing an order}

Fubini's theorem says that both orders work.  It does not say that both orders are equally
pleasant.

Consider
\[
    I=
    \int_0^1\int_0^1 2x e^{xy}\,dy\,dx.
\]
If we integrate with respect to \(y\) first, \(x\) is a constant:
\[
\begin{aligned}
    \int_0^1 2x e^{xy}\,dy
    &=
    2\int_0^1 x e^{xy}\,dy\\
    &=
    2\Bigl[e^{xy}\Bigr]_{y=0}^{y=1}\\
    &=
    2(e^x-1).
\end{aligned}
\]
Then
\[
\begin{aligned}
    I
    &=
    \int_0^1 2(e^x-1)\,dx\\
    &=
    2\Bigl[e^x-x\Bigr]_0^1\\
    &=
    2(e-1)-2\\
    &=
    2e-4.
\end{aligned}
\]
So
\[
\boxed{
    \int_0^1\int_0^1 2x e^{xy}\,dy\,dx=2e-4.
}
\]

The other order is possible, but less direct.  The moral is simple: before computing,
look for the order in which the inner integral has an easy antiderivative.

\subsection{Why the hypotheses matter}

Fubini's theorem is safe for continuous functions on rectangles.  If we sneak in an
unbounded singularity, the story can break.

Define
\[
    f(x,y)
    =
    \frac{x^2-y^2}{(x^2+y^2)^2}
\]
away from the origin in the square
\[
    0\le x\le1,
    \qquad
    0\le y\le1.
\]
This function is not continuous at \((0,0)\), and it is unbounded there.

If we read the iterated integrals as improper integrals, the two orders give different
answers.

For fixed \(x>0\),
\[
    \frac{\partial}{\partial y}
    \left(
        \frac{y}{x^2+y^2}
    \right)
    =
    \frac{x^2-y^2}{(x^2+y^2)^2}.
\]
Therefore
\[
\begin{aligned}
    \int_0^1
    \frac{x^2-y^2}{(x^2+y^2)^2}\,dy
    &=
    \left[
        \frac{y}{x^2+y^2}
    \right]_0^1\\
    &=
    \frac{1}{1+x^2}.
\end{aligned}
\]
Then
\[
    \int_0^1 \frac{1}{1+x^2}\,dx=\frac{\pi}{4}.
\]

But for fixed \(y>0\),
\[
    \frac{\partial}{\partial x}
    \left(
        -\frac{x}{x^2+y^2}
    \right)
    =
    \frac{x^2-y^2}{(x^2+y^2)^2}.
\]
So
\[
\begin{aligned}
    \int_0^1
    \frac{x^2-y^2}{(x^2+y^2)^2}\,dx
    &=
    \left[
        -\frac{x}{x^2+y^2}
    \right]_0^1\\
    &=
    -\frac{1}{1+y^2}.
\end{aligned}
\]
Then
\[
    \int_0^1 -\frac{1}{1+y^2}\,dy=-\frac{\pi}{4}.
\]

The two answers are
\[
    \frac{\pi}{4}
    \qquad\text{and}\qquad
    -\frac{\pi}{4}.
\]
This is not a contradiction.  Fubini's theorem did not apply.  The function was not
continuous on the rectangle, and the ordinary double integral was not a proper double
integral over the square.

The warning is not meant to make the theorem feel fragile.  It is meant to protect the
theorem from being used outside its home.  For continuous functions on rectangles, the
order may be chosen for convenience.

\subsection{What comes next}

Rectangles are clean.  Their bounds do not change:
\[
    a\le x\le b,
    \qquad
    c\le y\le d.
\]
But most regions are not rectangular.  A pond has a curved shore.  A triangular plate has
a slanted edge.  A disk has circular boundaries.

The next step is to let the limits of the inner integral depend on the outer variable.

That is how double integrals escape the rectangle.

\section{Double integrals over general regions}
\label{sec:10.4}

A rectangle is easy to slice because every vertical slice has the same \(y\)-limits, and
every horizontal slice has the same \(x\)-limits.  A general region asks for more care.

The main rule is not a formula.  It is a habit:

\[
\boxed{
    \text{Draw the region before writing the integral.}
}
\]

The picture tells us the limits.

\subsection{A triangular flower bed}

A triangular flower bed is bounded by
\[
    y=0,
    \qquad
    x=4,
    \qquad
    y=\frac{x}{2}.
\]
Suppose the height of a layer of mulch spread over the bed is modeled by
\[
    m(x,y)=6-y.
\]
Find the total volume of mulch over the bed.

The region is
\[
    D=\left\{(x,y):0\le x\le4,\ 0\le y\le \frac{x}{2}\right\}.
\]
For each fixed \(x\), the vertical slice runs from the bottom edge
\[
    y=0
\]
to the slanted edge
\[
    y=\frac{x}{2}.
\]

\begin{figure}[h]
\centering
\begin{tikzpicture}[scale=1.05, >=stealth]
    \draw[->] (-0.3,0) -- (5.0,0) node[right] {\(x\)};
    \draw[->] (0,-0.3) -- (0,3.0) node[above] {\(y\)};

    \draw[fill=green!12, draw=green!45!black, thick]
        (0,0) -- (4,0) -- (4,2) -- cycle;

    \draw[thick] (0,0) -- (4,2);
    
    % Shifted the formula to the right side of the line to avoid overlap
    \node[above left] at (3.6,1.8) {\(y=x/2\)};

    \draw[fill=blue!25, draw=blue!60!black, thick]
        (2.2,0) -- (2.55,0) -- (2.55,1.275) -- (2.2,1.1) -- cycle;

    \node[blue!60!black] at (1.45,1.7) {vertical slice};
    \draw[->, blue!60!black, thick] (1.75,1.55) -- (2.3,0.9);

    \node[below] at (4,0) {\(4\)};
    \node[left] at (0,2) {\(2\)};
    \fill (4,2) circle (2pt);
    \node[right] at (4,2) {\((4,2)\)};
\end{tikzpicture}
\caption{For this triangle, a vertical slice has variable height \(x/2\).}
\end{figure}

Therefore
\[
    V=\iint_D (6-y)\,dA
\]
becomes
\[
    V=
    \int_0^4\int_0^{x/2} (6-y)\,dy\,dx.
\]
Compute the inner integral:
\[
\begin{aligned}
    \int_0^{x/2} (6-y)\,dy
    &=
    \left[6y-\frac12y^2\right]_0^{x/2}\\
    &=
    3x-\frac{x^2}{8}.
\end{aligned}
\]
Now integrate in \(x\):
\[
\begin{aligned}
    V
    &=
    \int_0^4
    \left(
        3x-\frac{x^2}{8}
    \right)\,dx\\
    &=
    \left[
        \frac32x^2-\frac{x^3}{24}
    \right]_0^4\\
    &=
    24-\frac{64}{24}\\
    &=
    24-\frac{8}{3}\\
    &=
    \frac{64}{3}.
\end{aligned}
\]
Thus
\[
\boxed{
    V=\frac{64}{3}.
}
\]

The only new feature was the upper limit
\[
    \frac{x}{2}.
\]
The upper edge of the region became a limit of integration.

\subsection{Type I regions}

The triangular flower bed is an example of a region that is easy to slice vertically.

\begin{definition}[Type I region]
A region \(D\) in the plane is called a \emph{Type I region} if it can be written as
\[
\boxed{
    D=
    \left\{
        (x,y):
        a\le x\le b,\ 
        \phi_1(x)\le y\le \phi_2(x)
    \right\}.
}
\]
Here the \(x\)-values run from \(a\) to \(b\), and for each fixed \(x\), the \(y\)-values run
from a lower curve
\[
    y=\phi_1(x)
\]
to an upper curve
\[
    y=\phi_2(x).
\]
\end{definition}

For a Type I region,
\[
\boxed{
    \iint_D f(x,y)\,dA
    =
    \int_a^b
    \int_{\phi_1(x)}^{\phi_2(x)}
    f(x,y)\,dy\,dx.
}
\]

The inner integral moves along a vertical slice.  The outer integral moves those slices
from left to right.

\begin{figure}[h]
\centering
\begin{tikzpicture}[scale=1.0, >=stealth]
    \draw[->] (-0.4,0) -- (5.5,0) node[right] {\(x\)};
    \draw[->] (0,-0.5) -- (0,4.2) node[above] {\(y\)};

    \draw[fill=gray!10, draw=none]
        plot[domain=0.8:4.6, samples=80] (\x,{0.65+0.25*(\x-1.2)^2})
        -- plot[domain=4.6:0.8, samples=80] (\x,{3.1-0.18*(\x-2.8)^2})
        -- cycle;

    \draw[thick, blue!60!black, domain=0.8:4.6, samples=80]
        plot (\x,{0.65+0.25*(\x-1.2)^2});
    \draw[thick, red!65!black, domain=0.8:4.6, samples=80]
        plot (\x,{3.1-0.18*(\x-2.8)^2});

    \draw[dashed] (0.8,0) -- (0.8,{0.65+0.25*(0.8-1.2)^2});
    \draw[dashed] (4.6,0) -- (4.6,{0.65+0.25*(4.6-1.2)^2});

    \draw[very thick, green!50!black]
        (2.7,{0.65+0.25*(2.7-1.2)^2})
        -- (2.7,{3.1-0.18*(2.7-2.8)^2});

    \node[blue!60!black] at (2.0,0.65) {\(y=\phi_1(x)\)};
    \node[red!65!black] at (3.2,3.35) {\(y=\phi_2(x)\)};
    \node[below] at (0.8,0) {\(a\)};
    \node[below] at (4.6,0) {\(b\)};
    \node[green!50!black, right] at (2.75,2.0) {vertical slice};
\end{tikzpicture}
\caption{A Type I region is described by vertical slices.}
\end{figure}

\begin{theorem}[Double integral over a Type I region]
Suppose
\[
    D=
    \left\{
        (x,y):
        a\le x\le b,\ 
        \phi_1(x)\le y\le \phi_2(x)
    \right\},
\]
where \(\phi_1\) and \(\phi_2\) are continuous, and suppose \(f\) is continuous on \(D\).
Then
\[
\boxed{
    \iint_D f(x,y)\,dA
    =
    \int_a^b
    \int_{\phi_1(x)}^{\phi_2(x)}
    f(x,y)\,dy\,dx.
}
\]
\end{theorem}

\begin{proof}[Idea of the proof]
For each fixed \(x\), the vertical cross-section of \(D\) has \(y\)-values from
\[
    \phi_1(x)
\]
to
\[
    \phi_2(x).
\]
So the inner integral
\[
    \int_{\phi_1(x)}^{\phi_2(x)} f(x,y)\,dy
\]
adds \(f\) along that vertical slice.  The outer integral then adds all those slices from
\(x=a\) to \(x=b\).

This is the same slicing idea as Fubini's theorem, except now the slice endpoints can move
with \(x\).
\end{proof}

\subsection{Type II regions}

Some regions are easier to slice horizontally.

The same triangular flower bed can also be described by horizontal slices.  Since
\[
    y=\frac{x}{2}
\]
is the same as
\[
    x=2y,
\]
the triangle can be written as
\[
    D=
    \left\{
        (x,y):
        0\le y\le2,\ 
        2y\le x\le4
    \right\}.
\]
Thus
\[
\begin{aligned}
    V
    &=
    \int_0^2\int_{2y}^4 (6-y)\,dx\,dy\\
    &=
    \int_0^2 (6-y)(4-2y)\,dy\\
    &=
    \int_0^2 (24-16y+2y^2)\,dy\\
    &=
    \left[
        24y-8y^2+\frac23y^3
    \right]_0^2\\
    &=
    48-32+\frac{16}{3}\\
    &=
    \frac{64}{3}.
\end{aligned}
\]
This matches the answer from vertical slices:
\[
\boxed{
    \frac{64}{3}.
}
\]

\begin{definition}[Type II region]
A region \(D\) in the plane is called a \emph{Type II region} if it can be written as
\[
\boxed{
    D=
    \left\{
        (x,y):
        c\le y\le d,\ 
        \psi_1(y)\le x\le \psi_2(y)
    \right\}.
}
\]
Here the \(y\)-values run from \(c\) to \(d\), and for each fixed \(y\), the \(x\)-values run
from a left curve
\[
    x=\psi_1(y)
\]
to a right curve
\[
    x=\psi_2(y).
\]
\end{definition}

For a Type II region,
\[
\boxed{
    \iint_D f(x,y)\,dA
    =
    \int_c^d
    \int_{\psi_1(y)}^{\psi_2(y)}
    f(x,y)\,dx\,dy.
}
\]

\begin{figure}[h]
\centering
\begin{tikzpicture}[scale=1.0, >=stealth]
    \draw[->] (-0.4,0) -- (5.5,0) node[right] {\(x\)};
    \draw[->] (0,-0.5) -- (0,4.2) node[above] {\(y\)};

    \draw[fill=gray!10, draw=none]
        plot[domain=0.7:3.5, samples=80] ({0.7+0.18*(\x-2)^2},\x)
        -- plot[domain=3.5:0.7, samples=80] ({4.7-0.25*(\x-2.2)^2},\x)
        -- cycle;

    \draw[thick, blue!60!black, domain=0.7:3.5, samples=80]
        plot ({0.7+0.18*(\x-2)^2},\x);
    \draw[thick, red!65!black, domain=0.7:3.5, samples=80]
        plot ({4.7-0.25*(\x-2.2)^2},\x);

    \draw[dashed] (0,0.7) -- ({0.7+0.18*(0.7-2)^2},0.7);
    \draw[dashed] (0,3.5) -- ({0.7+0.18*(3.5-2)^2},3.5);

    \draw[very thick, green!50!black]
        ({0.7+0.18*(2.4-2)^2},2.4)
        -- ({4.7-0.25*(2.4-2.2)^2},2.4);

    \node[blue!60!black] at (1.45,1.1) {\(x=\psi_1(y)\)};
    \node[red!65!black] at (4.2,3.45) {\(x=\psi_2(y)\)};
    \node[left] at (0,0.7) {\(c\)};
    \node[left] at (0,3.5) {\(d\)};
    \node[green!50!black, above] at (3.0,2.45) {horizontal slice};
\end{tikzpicture}
\caption{A Type II region is described by horizontal slices.}
\end{figure}

\begin{theorem}[Double integral over a Type II region]
Suppose
\[
    D=
    \left\{
        (x,y):
        c\le y\le d,\ 
        \psi_1(y)\le x\le \psi_2(y)
    \right\},
\]
where \(\psi_1\) and \(\psi_2\) are continuous, and suppose \(f\) is continuous on \(D\).
Then
\[
\boxed{
    \iint_D f(x,y)\,dA
    =
    \int_c^d
    \int_{\psi_1(y)}^{\psi_2(y)}
    f(x,y)\,dx\,dy.
}
\]
\end{theorem}

\subsection{Drawing the region before integrating}

The hardest part of many double integral problems is not the integration.  It is setting
up the bounds.

A reliable routine is:

\begin{enumerate}
\item Draw the boundary curves.

\item Mark the intersection points.

\item Decide whether vertical slices or horizontal slices are simpler.

\item Write the outer limits from the shadow of the region on the outer axis.

\item Write the inner limits from the slice.
\end{enumerate}

The word \emph{shadow} is useful.  For a Type I setup, the region casts a shadow on the
\(x\)-axis:
\[
    a\le x\le b.
\]
For a Type II setup, it casts a shadow on the \(y\)-axis:
\[
    c\le y\le d.
\]

\subsection{Reversing the order of integration}

Sometimes one order leads to an integral we cannot evaluate with elementary functions.
The other order may be easy.

Consider
\[
    I=
    \int_0^1\int_y^1 e^{x^2}\,dx\,dy.
\]
The inner integral
\[
    \int_y^1 e^{x^2}\,dx
\]
has no elementary antiderivative.  But the region may save us.

The bounds say
\[
    0\le y\le1,
    \qquad
    y\le x\le1.
\]
So the region is the triangle
\[
    0\le y\le x\le1.
\]
Equivalently,
\[
    0\le x\le1,
    \qquad
    0\le y\le x.
\]

\begin{figure}[h]
\centering
\begin{tikzpicture}[scale=2.4, >=stealth]
    \draw[->] (-0.08,0) -- (1.25,0) node[right] {\(x\)};
    \draw[->] (0,-0.08) -- (0,1.25) node[above] {\(y\)};

    \draw[fill=purple!13, draw=purple!60!black, thick]
        (0,0) -- (1,0) -- (1,1) -- cycle;

    \draw[thick] (0,0) -- (1,1);
    \node[above left] at (0.52,0.52) {\(y=x\)};

    \draw[very thick, blue!60!black] (0.62,0) -- (0.62,0.62);
    \node[blue!60!black, right] at (0.63,0.33) {\(0\le y\le x\)};

    \node[below] at (1,0) {\(1\)};
    \node[left] at (0,1) {\(1\)};
    \fill (1,1) circle (0.8pt);
    \node[above right] at (1,1) {\((1,1)\)};
\end{tikzpicture}
\caption{The integral \(\int_0^1\int_y^1 e^{x^2}\,dx\,dy\) is over the triangle
\(0\le y\le x\le1\).}
\end{figure}

Now reverse the order:
\[
    I=
    \int_0^1\int_0^x e^{x^2}\,dy\,dx.
\]
The inner integral is now easy because \(e^{x^2}\) is constant with respect to \(y\):
\[
\begin{aligned}
    \int_0^x e^{x^2}\,dy
    &=
    xe^{x^2}.
\end{aligned}
\]
So
\[
\begin{aligned}
    I
    &=
    \int_0^1 xe^{x^2}\,dx.
\end{aligned}
\]
Let
\[
    u=x^2,
    \qquad
    du=2x\,dx.
\]
Then
\[
\begin{aligned}
    I
    &=
    \frac12\int_0^1 e^u\,du\\
    &=
    \frac12(e-1).
\end{aligned}
\]
Thus
\[
\boxed{
    \int_0^1\int_y^1 e^{x^2}\,dx\,dy
    =
    \frac{e-1}{2}.
}
\]

The integral became possible not because the function changed, but because the slices
changed.

\subsection{A region that must be split}

Not every region is Type I or Type II in one piece.  Sometimes a slice hits the region in
two separate intervals.  Then one pair of limits would fill in space that is not really part
of the region.

For example, let
\[
    D=
    \bigl([0,1]\times[0,1]\bigr)
    \cup
    \bigl([0,1]\times[2,3]\bigr).
\]
This is two square islands with a gap between them.

A vertical slice at a fixed \(x\) hits
\[
    0\le y\le1
\]
and also
\[
    2\le y\le3.
\]
It does not hit the gap
\[
    1<y<2.
\]
So writing
\[
    0\le x\le1,
    \qquad
    0\le y\le3
\]
would describe the full rectangle, not the two islands.

The correct setup is a sum:
\[
\boxed{
    \iint_D f(x,y)\,dA
    =
    \int_0^1\int_0^1 f(x,y)\,dy\,dx
    +
    \int_0^1\int_2^3 f(x,y)\,dy\,dx.
}
\]

The principle is simple:

\[
\boxed{
    \text{If one slice breaks into pieces, split the region.}
}
\]

\subsection{Regions with holes}

A hole is another reason to split or subtract.

Imagine an L-shaped floor:
\[
    D=
    [0,3]\times[0,2]
    \setminus
    [1,3]\times[1,2].
\]
It is a \(3\)-by-\(2\) rectangle with the upper-right rectangle removed.

Suppose
\[
    f(x,y)=1+x
\]
is the density of a thin coating on the floor.  The total amount of coating is
\[
    \iint_D (1+x)\,dA.
\]
One method is to split the L-shape into two rectangles:
\[
    D_1=[0,3]\times[0,1],
\]
and
\[
    D_2=[0,1]\times[1,2].
\]
Then
\[
\begin{aligned}
    \iint_D (1+x)\,dA
    &=
    \int_0^3\int_0^1 (1+x)\,dy\,dx
    +
    \int_0^1\int_1^2 (1+x)\,dy\,dx.
\end{aligned}
\]
Compute:
\[
\begin{aligned}
    \int_0^3\int_0^1 (1+x)\,dy\,dx
    &=
    \int_0^3 (1+x)\,dx\\
    &=
    \left[x+\frac12x^2\right]_0^3\\
    &=
    3+\frac92\\
    &=
    \frac{15}{2}.
\end{aligned}
\]
Also,
\[
\begin{aligned}
    \int_0^1\int_1^2 (1+x)\,dy\,dx
    &=
    \int_0^1 (1+x)\,dx\\
    &=
    \left[x+\frac12x^2\right]_0^1\\
    &=
    \frac32.
\end{aligned}
\]
Therefore
\[
\boxed{
    \iint_D (1+x)\,dA
    =
    \frac{15}{2}+\frac32
    =
    9.
}
\]

\begin{figure}[h]
\centering
\begin{tikzpicture}[scale=1.1, >=stealth]
    \draw[->] (-0.3,0) -- (4.0,0) node[right] {\(x\)};
    \draw[->] (0,-0.3) -- (0,3.0) node[above] {\(y\)};

    % The L-shaped region
    \draw[fill=orange!18, draw=orange!60!black, thick]
        (0,0) -- (3,0) -- (3,1) -- (1,1) -- (1,2) -- (0,2) -- cycle;

    % Line splitting D1 and D2
    \draw[dashed, thick] (0,1) -- (1,1);

    % Dashed outline showing the missing part of the full 3x2 rectangle
    \draw[dashed, gray!80, thick] (3,1) -- (3,2) -- (1,2);

    \node at (1.5,0.5) {\(D_1\)};
    \node at (0.5,1.5) {\(D_2\)};
    \node[gray!80!black] at (2.0,1.5) {hole};

    \node[below] at (1,0) {\(1\)};
    \node[below] at (3,0) {\(3\)};
    \node[left] at (0,1) {\(1\)};
    \node[left] at (0,2) {\(2\)};
\end{tikzpicture}
\caption{An L-shaped region is easiest to integrate over by splitting it into pieces.}
\end{figure}

Another method is subtraction:
\[
\iint_D (1+x)\,dA
=
\iint_{[0,3]\times[0,2]}(1+x)\,dA
-
\iint_{[1,3]\times[1,2]}(1+x)\,dA.
\]
The full rectangle contributes
\[
    \int_0^3\int_0^2(1+x)\,dy\,dx
    =
    2\int_0^3(1+x)\,dx
    =
    15.
\]
The removed rectangle contributes
\[
    \int_1^3\int_1^2(1+x)\,dy\,dx
    =
    \int_1^3(1+x)\,dx
    =
    6.
\]
So again
\[
    15-6=9.
\]

Both methods are correct.  The good method is the one that makes the region easiest to
describe.

\subsection{The main lesson}

For rectangles, the limits are constant:
\[
    a\le x\le b,
    \qquad
    c\le y\le d.
\]
For general regions, one pair of limits often depends on the other variable:
\[
    \phi_1(x)\le y\le\phi_2(x)
\]
or
\[
    \psi_1(y)\le x\le\psi_2(y).
\]

So the calculation is still based on one idea:

\[
\boxed{
    \text{add slices.}
}
\]

The only difference is that the slices may grow, shrink, split, or avoid holes.

The next obstacle is circular geometry.  A disk can be described with \(x\)- and
\(y\)-bounds, but the limits contain square roots.  A circular ring is even less friendly in
rectangular coordinates.

For circular regions, the natural slices are not vertical or horizontal.  They are radial
and angular.

That leads to polar coordinates.
% =====================================================================
%  Chapter 10: Double Integrals
%  Sections 10.5 and 10.6
% =====================================================================

\section{Double integrals in polar coordinates}
\label{sec:10.5}

Section \ref{sec:10.4} taught us how to integrate over regions that are not rectangles.
The main tool was slicing.  But circular regions still felt awkward in rectangular
coordinates.  A disk can be described by
\[
    -a\le x\le a,
    \qquad
    -\sqrt{a^2-x^2}\le y\le \sqrt{a^2-x^2},
\]
but those square roots are a sign that the coordinates are fighting the geometry.

For circular regions, the natural coordinates are
\[
    x=r\cos\theta,
    \qquad
    y=r\sin\theta.
\]
The new question is not just how to rewrite the region.  We must also understand how
area changes.

\subsection{A circular bowl}

A shallow circular bowl has height
\[
    z=9-x^2-y^2
\]
above the circular floor
\[
    x^2+y^2\le 9.
\]
The volume under the surface is
\[
    V=\iint_D (9-x^2-y^2)\,dA,
\]
where
\[
    D=\{(x,y):x^2+y^2\le9\}.
\]

In rectangular coordinates, the same integral becomes
\[
    \int_{-3}^{3}
    \int_{-\sqrt{9-x^2}}^{\sqrt{9-x^2}}
    (9-x^2-y^2)\,dy\,dx.
\]
This is correct, but it does not look like the bowl.

In polar coordinates,
\[
    x^2+y^2=r^2.
\]
The disk is simply
\[
    0\le r\le3,
    \qquad
    0\le\theta\le2\pi.
\]
The height becomes
\[
    9-x^2-y^2=9-r^2.
\]

But there is one more piece.  A tiny polar box does not have area
\[
    \Delta r\,\Delta\theta.
\]
It has area approximately
\[
    r\,\Delta r\,\Delta\theta.
\]
So the integral should be
\[
    V=
    \int_0^{2\pi}\int_0^3 (9-r^2)r\,dr\,d\theta.
\]
Compute:
\[
\begin{aligned}
    V
    &=
    \int_0^{2\pi}
    \left[
        \frac92r^2-\frac14r^4
    \right]_0^3
    d\theta \\
    &=
    \int_0^{2\pi}
    \left(
        \frac92\cdot9-\frac14\cdot81
    \right)
    d\theta\\
    &=
    \int_0^{2\pi}
    \frac{81}{4}\,d\theta\\
    &=
    \frac{81\pi}{2}.
\end{aligned}
\]
Thus
\[
\boxed{
    V=\frac{81\pi}{2}.
}
\]

The factor \(r\) is the new feature.  Everything in polar integration depends on it.

\subsection{Why the area element is \texorpdfstring{\(r\,dr\,d\theta\)}{r dr dtheta}}

A polar grid is made from circles and rays.  If we fix \(r\) and let \(\theta\) change, we
move around a circle.  If we fix \(\theta\) and let \(r\) change, we move along a ray.

A small polar rectangle has two side lengths:

\[
    \text{radial side length}\approx \Delta r,
\]
and
\[
    \text{circular side length}\approx r\,\Delta\theta.
\]
The second formula is arc length:
\[
    \text{arc length}=(\text{radius})(\text{angle}).
\]
Therefore
\[
    \Delta A
    \approx
    (r\,\Delta\theta)(\Delta r)
    =
    r\,\Delta r\,\Delta\theta.
\]

\begin{figure}[h]
\centering
\begin{tikzpicture}[scale=1.15, >=stealth]
    \draw[->] (-0.3,0) -- (4.8,0) node[right] {\(x\)};
    \draw[->] (0,-0.3) -- (0,4.2) node[above] {\(y\)};

    \draw[gray!25] (0,0) circle (1.4);
    \draw[gray!25] (0,0) circle (2.2);
    \draw[gray!25] (0,0) circle (3.0);

    \draw[gray!40] (0,0) -- (4.0,0.0);
    \draw[gray!40] (0,0) -- ({4*cos(20)},{4*sin(20)});
    \draw[gray!40] (0,0) -- ({4*cos(40)},{4*sin(40)});
    \draw[gray!40] (0,0) -- ({4*cos(60)},{4*sin(60)});

    \filldraw[fill=blue!20, draw=blue!60!black, thick]
        ({2.2*cos(20)},{2.2*sin(20)})
        arc[start angle=20,end angle=40,radius=2.2]
        -- ({3.0*cos(40)},{3.0*sin(40)})
        arc[start angle=40,end angle=20,radius=3.0]
        -- cycle;

    \draw[<->, blue!70!black]
        ({2.45*cos(20)},{2.45*sin(20)})
        -- ({2.95*cos(20)},{2.95*sin(20)})
        node[midway, below right] {\(\Delta r\)};

    \node[blue!70!black] at ({3.35*cos(30)},{3.35*sin(30)})
        {\(r\,\Delta\theta\)};

    % Inner arc placed at r=0.8, label placed safely inside at r=0.55
    \draw[->, thick] (0.8,0) arc[start angle=0,end angle=20,radius=0.8];
    \node at ({0.55*cos(10)},{0.55*sin(10)}) {\(\theta_1\)};
    
    % Outer arc placed at r=1.15 to clear the inner label, label outside at r=1.4
    \draw[->, thick] (1.15,0) arc[start angle=0,end angle=40,radius=1.15];
    \node at ({1.4*cos(35)},{1.4*sin(35)}) {\(\theta_2\)};

    \node at (2.2,-0.55) {polar grid cells grow as \(r\) grows};
\end{tikzpicture}
\caption{A small polar rectangle has area approximately \(r\,\Delta r\,\Delta\theta\).}
\end{figure}

There is also a differential-forms way to see the same factor.

From
\[
    x=r\cos\theta,
    \qquad
    y=r\sin\theta,
\]
we get
\[
    dx=\cos\theta\,dr-r\sin\theta\,d\theta,
\]
and
\[
    dy=\sin\theta\,dr+r\cos\theta\,d\theta.
\]
Now wedge:
\[
\begin{aligned}
    dx\wedge dy
    &=
    (\cos\theta\,dr-r\sin\theta\,d\theta)
    \wedge
    (\sin\theta\,dr+r\cos\theta\,d\theta)\\
    &=
    r\cos^2\theta\,dr\wedge d\theta
    +
    r\sin^2\theta\,dr\wedge d\theta\\
    &=
    r(\cos^2\theta+\sin^2\theta)\,dr\wedge d\theta\\
    &=
    r\,dr\wedge d\theta.
\end{aligned}
\]
So the standard area form becomes
\[
\boxed{
    dx\wedge dy=r\,dr\wedge d\theta.
}
\]
For ordinary positive area,
\[
\boxed{
    dA=r\,dr\,d\theta.
}
\]

This is the same factor we saw geometrically.

\subsection{Polar rectangles}

A rectangle in the \((r,\theta)\)-plane becomes a curved region in the \(xy\)-plane.

\begin{definition}[Polar rectangle]
A \emph{polar rectangle} is a region described by
\[
\boxed{
    a\le r\le b,
    \qquad
    \alpha\le\theta\le\beta.
}
\]
In the \(xy\)-plane, this is the region between the circles
\[
    r=a
    \qquad\text{and}\qquad
    r=b
\]
and between the rays
\[
    \theta=\alpha
    \qquad\text{and}\qquad
    \theta=\beta.
\]
\end{definition}

For example,
\[
    1\le r\le3,
    \qquad
    0\le\theta\le\pi
\]
is a half-annulus.

Its area is
\[
\begin{aligned}
    A
    &=
    \int_0^\pi\int_1^3 r\,dr\,d\theta\\
    &=
    \int_0^\pi
    \left[
        \frac12r^2
    \right]_1^3
    d\theta\\
    &=
    \int_0^\pi 4\,d\theta\\
    &=
    4\pi.
\end{aligned}
\]
This agrees with geometry:
\[
    \frac12\pi(3^2)-\frac12\pi(1^2)=\frac92\pi-\frac12\pi=4\pi.
\]

\begin{figure}[h]
\centering
\begin{tikzpicture}[scale=0.95, >=stealth]
    \draw[->] (-3.8,0) -- (3.8,0) node[right] {\(x\)};
    \draw[->] (0,-0.5) -- (0,3.9) node[above] {\(y\)};

    \fill[orange!18]
        (1,0) arc[start angle=0,end angle=180,radius=1]
        -- (-3,0) arc[start angle=180,end angle=0,radius=3]
        -- cycle;

    \draw[orange!60!black, thick] (1,0) arc[start angle=0,end angle=180,radius=1];
    \draw[orange!60!black, thick] (3,0) arc[start angle=0,end angle=180,radius=3];
    \draw[orange!60!black, thick] (1,0) -- (3,0);
    \draw[orange!60!black, thick] (-1,0) -- (-3,0);

    \node at (0,2.1) {\(1\le r\le3\)};
    \node at (0,1.55) {\(0\le\theta\le\pi\)};
\end{tikzpicture}
\caption{A polar rectangle can be a curved sector or annular sector.}
\end{figure}

\subsection{The polar-coordinate formula}

Now we can state the general rule.

\begin{theorem}[Double integrals in polar coordinates]
Suppose \(D\) is described in polar coordinates by
\[
    D=
    \left\{
        (r,\theta):
        \alpha\le\theta\le\beta,\ 
        g(\theta)\le r\le h(\theta)
    \right\},
\]
where \(g\) and \(h\) are continuous, \(0\le g(\theta)\le h(\theta)\), and the polar
description covers the region once except possibly along boundary curves.

If \(f\) is continuous on \(D\), then
\[
\boxed{
    \iint_D f(x,y)\,dA
    =
    \int_\alpha^\beta
    \int_{g(\theta)}^{h(\theta)}
    f(r\cos\theta,r\sin\theta)\,r\,dr\,d\theta.
}
\]
Equivalently, in form language,
\[
\boxed{
    \int_D f(x,y)\,dx\wedge dy
    =
    \int_\alpha^\beta
    \int_{g(\theta)}^{h(\theta)}
    f(r\cos\theta,r\sin\theta)\,r\,dr\,d\theta.
}
\]
\end{theorem}

\begin{proof}[Idea of the proof]
Divide the region into small polar rectangles.  A small piece at radius \(r\) has area
approximately
\[
    r\,\Delta r\,\Delta\theta.
\]
On that piece,
\[
    f(x,y)=f(r\cos\theta,r\sin\theta).
\]
So a Riemann sum looks like
\[
    \sum
    f(r\cos\theta,r\sin\theta)\,r\,\Delta r\,\Delta\theta.
\]
As the polar pieces get small, this sum becomes
\[
    \int_\alpha^\beta
    \int_{g(\theta)}^{h(\theta)}
    f(r\cos\theta,r\sin\theta)\,r\,dr\,d\theta.
\]
The same conclusion follows from
\[
    dx\wedge dy=r\,dr\wedge d\theta.
\]
\end{proof}

There are two warnings hidden in the theorem.

First, the function \(f\) is assumed continuous.  This is a safe hypothesis that guarantees
the usual Riemann integral exists.  A wildly discontinuous function, such as the function
that is \(1\) at rational-coordinate points and \(0\) at irrational-coordinate points inside
a disk, is not Riemann integrable.  Changing coordinates cannot fix a function whose
ordinary integral is not defined.

Second, the polar description must not cover the region more than once.

For example,
\[
    0\le r\le1,
    \qquad
    0\le\theta\le4\pi
\]
traces the unit disk twice.  If we integrate \(1\), we get
\[
\begin{aligned}
    \int_0^{4\pi}\int_0^1 r\,dr\,d\theta
    &=
    \int_0^{4\pi}\frac12\,d\theta\\
    &=
    2\pi.
\end{aligned}
\]
But the area of the unit disk is only
\[
    \pi.
\]
The formula did not fail.  The setup counted every point twice.

\[
\boxed{
    \text{A repeated polar description gives a repeated integral.}
}
\]

\subsection{Radial symmetry}

A function is radially symmetric if it depends only on the distance from the origin:
\[
    f(x,y)=F(r),
    \qquad
    r=\sqrt{x^2+y^2}.
\]
On a disk or annulus centered at the origin, polar coordinates make this especially
simple.

If
\[
    D=\{(x,y):0\le r\le a\},
\]
then
\[
\begin{aligned}
    \iint_D F(r)\,dA
    &=
    \int_0^{2\pi}\int_0^a F(r)\,r\,dr\,d\theta\\
    &=
    2\pi\int_0^a F(r)\,r\,dr.
\end{aligned}
\]
So for radially symmetric functions on centered disks,
\[
\boxed{
    \iint_D F(r)\,dA
    =
    2\pi\int_0^a F(r)\,r\,dr.
}
\]

The factor \(2\pi\) comes from going all the way around the circle.  The factor \(r\)
still comes from area.

\begin{example}[A radially symmetric mound]
A circular pile of sand has height
\[
    h(r)=12-3r
\]
over the disk
\[
    0\le r\le4.
\]
The total volume is
\[
\begin{aligned}
    V
    &=
    \int_0^{2\pi}\int_0^4 (12-3r)r\,dr\,d\theta\\
    &=
    2\pi\int_0^4 (12r-3r^2)\,dr\\
    &=
    2\pi\left[6r^2-r^3\right]_0^4\\
    &=
    2\pi(96-64)\\
    &=
    64\pi.
\end{aligned}
\]
Thus
\[
\boxed{
    V=64\pi.
}
\]
\end{example}

\subsection{A shifted circle}

Polar coordinates are not only for regions centered at the origin.  They are often useful
whenever a boundary has a simple polar equation.

Consider the disk
\[
    D:\quad x^2+y^2\le2x.
\]
Complete the square:
\[
    x^2-2x+y^2\le0,
\]
so
\[
    (x-1)^2+y^2\le1.
\]
This is the disk of radius \(1\) centered at \((1,0)\).

In polar coordinates,
\[
    x^2+y^2=r^2,
    \qquad
    x=r\cos\theta.
\]
So the boundary
\[
    x^2+y^2=2x
\]
becomes
\[
    r^2=2r\cos\theta.
\]
For \(r\ge0\), this gives
\[
    r=2\cos\theta.
\]
This is nonnegative when
\[
    -\frac{\pi}{2}\le\theta\le\frac{\pi}{2}.
\]
Thus the disk is
\[
    -\frac{\pi}{2}\le\theta\le\frac{\pi}{2},
    \qquad
    0\le r\le2\cos\theta.
\]

\begin{figure}[h]
\centering
\begin{tikzpicture}[scale=1.05, >=stealth]
    \draw[->] (-0.8,0) -- (3.2,0) node[right] {\(x\)};
    \draw[->] (0,-1.7) -- (0,1.7) node[above] {\(y\)};

    \draw[fill=blue!12, draw=blue!60!black, thick] (1,0) circle (1);

    \fill (1,0) circle (2pt);
    \node[below] at (1,0) {\((1,0)\)};

    \draw[thick, red!65!black] (0,0) -- ({2*cos(35)*cos(35)},{2*cos(35)*sin(35)});
    \node[red!65!black, right] at (1.25,0.55) {\(r=2\cos\theta\)};

    \draw[->] (0.55,0) arc[start angle=0,end angle=35,radius=0.55];
    \node at (0.72,0.18) {\(\theta\)};

    \node at (1,-1.35) {\((x-1)^2+y^2\le1\)};
\end{tikzpicture}
\caption{The shifted disk \((x-1)^2+y^2\le1\) has polar boundary \(r=2\cos\theta\).}
\end{figure}

Now find
\[
    \iint_D (x^2+y^2)\,dA.
\]
Since
\[
    x^2+y^2=r^2,
\]
we get
\[
\begin{aligned}
    \iint_D (x^2+y^2)\,dA
    &=
    \int_{-\pi/2}^{\pi/2}
    \int_0^{2\cos\theta}
    r^2\,r\,dr\,d\theta\\
    &=
    \int_{-\pi/2}^{\pi/2}
    \left[
        \frac14r^4
    \right]_0^{2\cos\theta}
    d\theta\\
    &=
    \int_{-\pi/2}^{\pi/2}
    4\cos^4\theta\,d\theta.
\end{aligned}
\]
Using
\[
    \cos^4\theta
    =
    \left(\frac{1+\cos2\theta}{2}\right)^2
    =
    \frac38+\frac12\cos2\theta+\frac18\cos4\theta,
\]
we get
\[
\begin{aligned}
    \iint_D (x^2+y^2)\,dA
    &=
    4\left[
        \frac38\theta
        +\frac14\sin2\theta
        +\frac1{32}\sin4\theta
    \right]_{-\pi/2}^{\pi/2}\\
    &=
    4\cdot\frac{3\pi}{8}\\
    &=
    \frac{3\pi}{2}.
\end{aligned}
\]
Thus
\[
\boxed{
    \iint_D (x^2+y^2)\,dA=\frac{3\pi}{2}.
}
\]

\subsection{A spiral region}

Polar coordinates also describe regions bounded by spirals.

Suppose a garden path has boundary
\[
    r=\frac{\theta}{2\pi},
    \qquad
    0\le\theta\le2\pi.
\]
The region inside one turn of the spiral is
\[
    0\le\theta\le2\pi,
    \qquad
    0\le r\le \frac{\theta}{2\pi}.
\]
Its area is
\[
\begin{aligned}
    A
    &=
    \int_0^{2\pi}
    \int_0^{\theta/(2\pi)}
    r\,dr\,d\theta\\
    &=
    \int_0^{2\pi}
    \frac12\left(\frac{\theta}{2\pi}\right)^2
    d\theta\\
    &=
    \frac{1}{8\pi^2}
    \left[
        \frac{\theta^3}{3}
    \right]_0^{2\pi}\\
    &=
    \frac{1}{8\pi^2}\cdot\frac{8\pi^3}{3}\\
    &=
    \frac{\pi}{3}.
\end{aligned}
\]
So
\[
\boxed{
    A=\frac{\pi}{3}.
}
\]

\begin{figure}[h]
\centering
\begin{tikzpicture}[scale=2.1, >=stealth]
    \draw[->] (-1.25,0) -- (1.35,0) node[right] {\(x\)};
    \draw[->] (0,-1.25) -- (0,1.35) node[above] {\(y\)};

    \draw[fill=green!13, draw=none]
        plot[domain=0:6.283, samples=200, smooth, variable=\t]
        ({(\t/(2*pi))*cos(\t r)}, {(\t/(2*pi))*sin(\t r)})
        -- (0,0)
        -- cycle;

    \draw[thick, green!45!black, domain=0:6.283, samples=200, smooth, variable=\t]
        plot ({(\t/(2*pi))*cos(\t r)}, {(\t/(2*pi))*sin(\t r)});

    \node[green!45!black] at (0.15,-0.8) {\(r=\theta/(2\pi)\)};
    \node at (0,-1.45) {area \(=\pi/3\)};
\end{tikzpicture}
\caption{A spiral region can have very simple polar bounds.}
\end{figure}

The rectangular description of the same region would be unpleasant.  The polar
description is almost a sentence:
\[
    \text{as \(\theta\) turns once, \(r\) grows from \(0\) to \(1\).}
\]

\subsection{The main lesson}

Changing to polar coordinates does two things:

\[
    x=r\cos\theta,\qquad y=r\sin\theta,
\]
and
\[
    dA=r\,dr\,d\theta.
\]

The first line changes the function and the region.  The second line changes the area
element.

So the safe routine is:

\begin{enumerate}
\item Draw the region.

\item Rewrite the boundary in polar form.

\item Rewrite the integrand using
\[
    x=r\cos\theta,\qquad y=r\sin\theta.
\]

\item Multiply by the area factor \(r\).

\item Check that the polar description covers the region once.
\end{enumerate}

The factor \(r\) is not optional.  It is the price of using circular coordinates.

Now that we can integrate over many regions, we can ask what double integrals are
good for besides volume.  The same symbol
\[
    \iint_D f\,dA
\]
can measure mass, center of mass, rotational inertia, probability, average temperature,
rainfall, or elevation.

The integral is the same.  The meaning of \(f\) changes.

\section{Applications of double integrals}
\label{sec:10.6}

A double integral adds tiny pieces:
\[
    f(x,y)\,\Delta A.
\]
If \(f\) is height, the pieces are tiny volumes.  But \(f\) does not have to be height.
It can be density, temperature, probability, or any quantity spread over a region.

This is why double integrals are useful.  They turn a local measurement into a total
quantity.

\subsection{Mass from density}

A thin metal sign occupies the rectangle
\[
    D=[0,2]\times[0,1].
\]
Because of a coating process, the sign is heavier toward the east.  Its density is
\[
    \sigma(x,y)=1+x
\]
kilograms per square meter.

A tiny piece of area \(\Delta A\) near \((x,y)\) has approximate mass
\[
    \sigma(x,y)\,\Delta A.
\]
So the total mass should be
\[
    m=\iint_D \sigma(x,y)\,dA.
\]
For this sign,
\[
\begin{aligned}
    m
    &=
    \int_0^2\int_0^1 (1+x)\,dy\,dx\\
    &=
    \int_0^2 (1+x)\,dx\\
    &=
    \left[
        x+\frac12x^2
    \right]_0^2\\
    &=
    2+2\\
    &=
    4.
\end{aligned}
\]
Thus the sign has mass
\[
\boxed{
    m=4\text{ kg}.
}
\]

\begin{definition}[Mass of a lamina]
A thin plate in the plane is often called a \emph{lamina}.  If a lamina occupies a region
\(D\) and has density
\[
    \sigma(x,y)
\]
measured in mass per unit area, then its total mass is
\[
\boxed{
    m=\iint_D \sigma(x,y)\,dA.
}
\]
In form language,
\[
\boxed{
    m=\int_D \sigma(x,y)\,dx\wedge dy.
}
\]
\end{definition}

Density must be nonnegative for a physical mass distribution:
\[
    \sigma(x,y)\ge0.
\]
Also, the total mass must be positive if we want to speak about a center of mass.

A signed function can still be integrated, but then it describes something like signed
charge, not ordinary mass.  For instance, on the symmetric rectangle
\[
    [-1,1]\times[0,1],
\]
the signed density
\[
    \sigma(x,y)=x
\]
has total integral
\[
    \iint_D x\,dA=0.
\]
That may make sense as total charge.  It does not make sense as a mass whose center we
can find by dividing by total mass.

\subsection{Center of mass}

The same coated sign from above has density
\[
    \sigma(x,y)=1+x
\]
on
\[
    D=[0,2]\times[0,1].
\]
Its mass is
\[
    m=4.
\]

Where should we balance it?

For a collection of point masses, the center of mass is a weighted average of positions.
For a continuous lamina, the sums become integrals.

The \(x\)-coordinate of the center of mass should be the average of \(x\), weighted by
mass:
\[
    \bar x
    =
    \frac{1}{m}\iint_D x\,\sigma(x,y)\,dA.
\]
The \(y\)-coordinate is
\[
    \bar y
    =
    \frac{1}{m}\iint_D y\,\sigma(x,y)\,dA.
\]

For our sign,
\[
\begin{aligned}
    \iint_D x\sigma(x,y)\,dA
    &=
    \int_0^2\int_0^1 x(1+x)\,dy\,dx\\
    &=
    \int_0^2 (x+x^2)\,dx\\
    &=
    \left[
        \frac12x^2+\frac13x^3
    \right]_0^2\\
    &=
    2+\frac83\\
    &=
    \frac{14}{3}.
\end{aligned}
\]
So
\[
    \bar x=\frac{1}{4}\cdot\frac{14}{3}=\frac76.
\]

Similarly,
\[
\begin{aligned}
    \iint_D y\sigma(x,y)\,dA
    &=
    \int_0^2\int_0^1 y(1+x)\,dy\,dx\\
    &=
    \int_0^2
    (1+x)
    \left[
        \frac12y^2
    \right]_0^1
    dx\\
    &=
    \frac12\int_0^2(1+x)\,dx\\
    &=
    \frac12\cdot4\\
    &=
    2.
\end{aligned}
\]
Thus
\[
    \bar y=\frac{2}{4}=\frac12.
\]
The center of mass is
\[
\boxed{
    \left(\bar x,\bar y\right)=\left(\frac76,\frac12\right).
}
\]

The answer makes sense.  If the sign had uniform density, the center would be
\[
    (1,1/2).
\]
Since the density is larger when \(x\) is larger, the center moves right to
\[
    \bar x=\frac76>1.
\]
But the density does not depend on \(y\), so the vertical coordinate stays
\[
    \bar y=\frac12.
\]

\begin{figure}[h]
\centering
\begin{tikzpicture}[scale=1.4, >=stealth]
    \draw[->] (-0.25,0) -- (2.65,0) node[right] {\(x\)};
    \draw[->] (0,-0.25) -- (0,1.55) node[above] {\(y\)};

    % Replaced uniform fill with a gradient to show increasing density
    \draw[left color=gray!5, right color=gray!45, draw=black, thick] (0,0) rectangle (2,1);

    \foreach \x in {0.25,0.55,0.85,1.15,1.45,1.75}
    {
        \pgfmathsetmacro{\rad}{0.7+0.35*\x}
        \fill[blue!60] (\x,0.25) circle ({0.025*\rad});
        \fill[blue!60] (\x,0.50) circle ({0.025*\rad});
        \fill[blue!60] (\x,0.75) circle ({0.025*\rad});
    }

    \fill[red!70!black] (1.1667,0.5) circle (2pt);
    \draw[red!70!black, dashed] (1.1667,0) -- (1.1667,0.5);
    \draw[red!70!black, dashed] (0,0.5) -- (1.1667,0.5);
    \node[red!70!black, above right] at (1.1667,0.5)
        {\(\left(\frac76,\frac12\right)\)};

    \node[below] at (2,0) {\(2\)};
    \node[left] at (0,1) {\(1\)};
    \node at (1,-0.55) {density \(\sigma(x,y)=1+x\)};
\end{tikzpicture}
\caption{The density increases to the right, so the center of mass shifts right.}
\end{figure}

\begin{definition}[Moments and center of mass]
Let a lamina occupy a region \(D\), with density \(\sigma(x,y)\), and total mass
\[
    m=\iint_D \sigma(x,y)\,dA>0.
\]
The \emph{moment about the \(y\)-axis} is
\[
\boxed{
    M_y=\iint_D x\,\sigma(x,y)\,dA.
}
\]
The \emph{moment about the \(x\)-axis} is
\[
\boxed{
    M_x=\iint_D y\,\sigma(x,y)\,dA.
}
\]
The center of mass is
\[
\boxed{
    \bar x=\frac{M_y}{m},
    \qquad
    \bar y=\frac{M_x}{m}.
}
\]
\end{definition}

The notation can feel backwards at first:
\[
    M_y
\]
uses \(x\), because \(x\) measures distance from the \(y\)-axis.  Similarly,
\[
    M_x
\]
uses \(y\), because \(y\) measures distance from the \(x\)-axis.

If the density is constant, the center of mass is the geometric center of the region,
called the \emph{centroid}.

\subsection{Moment of inertia}

Mass tells us how hard it is to move an object in a straight line.  Moment of inertia
tells us how hard it is to rotate an object.

A small mass far from the axis of rotation matters more than the same mass near the
axis.  The distance is squared.

Consider a uniform circular disk of radius \(a\), centered at the origin, with constant
density
\[
    \sigma_0.
\]
Find its moment of inertia about the origin, as if it were spinning around an axle through
the center perpendicular to the disk.

A point at distance \(r\) from the origin contributes a factor
\[
    r^2.
\]
The density is \(\sigma_0\), and
\[
    dA=r\,dr\,d\theta.
\]
So
\[
\begin{aligned}
    I_0
    &=
    \int_0^{2\pi}\int_0^a r^2\sigma_0\,r\,dr\,d\theta\\
    &=
    \sigma_0\int_0^{2\pi}\int_0^a r^3\,dr\,d\theta\\
    &=
    \sigma_0(2\pi)
    \left[
        \frac14r^4
    \right]_0^a\\
    &=
    \frac{\pi\sigma_0a^4}{2}.
\end{aligned}
\]
The mass of the disk is
\[
    m=\sigma_0\pi a^2.
\]
Therefore
\[
\boxed{
    I_0=\frac12ma^2.
}
\]

This formula says that increasing the radius has a strong effect.  The mass grows like
\(a^2\), but the moment of inertia about the center grows like \(a^4\) if density stays
fixed.

\begin{definition}[Moments of inertia]
Let a lamina occupy a region \(D\), with density \(\sigma(x,y)\).

The moment of inertia about the \(x\)-axis is
\[
\boxed{
    I_x=\iint_D y^2\sigma(x,y)\,dA.
}
\]
The moment of inertia about the \(y\)-axis is
\[
\boxed{
    I_y=\iint_D x^2\sigma(x,y)\,dA.
}
\]
The moment of inertia about the origin is
\[
\boxed{
    I_0=\iint_D (x^2+y^2)\sigma(x,y)\,dA.
}
\]
Since
\[
    x^2+y^2
\]
is the squared distance from the origin,
\[
\boxed{
    I_0=I_x+I_y.
}
\]
\end{definition}

In polar coordinates,
\[
    x^2+y^2=r^2,
\]
so moments of inertia about the origin often become especially simple:
\[
    I_0=\iint_D r^2\sigma\,dA.
\]

\subsection{Probability density over a region}

A probability density is another kind of density.  Instead of mass per unit area, it
measures probability per unit area.

Suppose a dart lands in the unit disk
\[
    0\le r\le1,
    \qquad
    0\le\theta\le2\pi.
\]
The dart is more likely to land near the center than near the edge.  A simple model is
\[
    p(r)=C(1-r),
\]
where \(C\) is a constant.

To be a probability density, the total probability must be \(1\):
\[
    \iint_D p\,dA=1.
\]
So
\[
\begin{aligned}
    1
    &=
    \int_0^{2\pi}\int_0^1 C(1-r)\,r\,dr\,d\theta\\
    &=
    2\pi C\int_0^1(r-r^2)\,dr\\
    &=
    2\pi C
    \left[
        \frac12r^2-\frac13r^3
    \right]_0^1\\
    &=
    2\pi C\left(\frac12-\frac13\right)\\
    &=
    \frac{\pi C}{3}.
\end{aligned}
\]
Thus
\[
    C=\frac{3}{\pi}.
\]
The density is
\[
    p(r)=\frac{3}{\pi}(1-r).
\]

Now find the probability that the dart lands within half the radius:
\[
    0\le r\le\frac12.
\]
Compute:
\[
\begin{aligned}
    P\left(r\le\frac12\right)
    &=
    \int_0^{2\pi}
    \int_0^{1/2}
    \frac{3}{\pi}(1-r)r\,dr\,d\theta\\
    &=
    6\int_0^{1/2}(r-r^2)\,dr\\
    &=
    6\left[
        \frac12r^2-\frac13r^3
    \right]_0^{1/2}\\
    &=
    6\left(
        \frac18-\frac1{24}
    \right)\\
    &=
    6\cdot\frac{2}{24}\\
    &=
    \frac12.
\end{aligned}
\]
So
\[
\boxed{
    P\left(r\le\frac12\right)=\frac12.
}
\]

For a uniform dartboard, the probability inside half the radius would be
\[
    \frac{\pi(1/2)^2}{\pi(1)^2}=\frac14.
\]
Here it is larger because the density is concentrated near the center.

\begin{definition}[Probability density over a region]
A function \(p(x,y)\) is a probability density on a region \(D\) if
\[
\boxed{
    p(x,y)\ge0
}
\]
and
\[
\boxed{
    \iint_D p(x,y)\,dA=1.
}
\]
For a subregion \(A\subseteq D\),
\[
\boxed{
    P((X,Y)\in A)=\iint_A p(x,y)\,dA.
}
\]
\end{definition}

Both conditions matter.

If
\[
    p(x,y)=2
\]
on the unit square, then
\[
    \iint_{[0,1]\times[0,1]}2\,dA=2,
\]
so it cannot be a probability density.  Total probability would be \(2\), which is
impossible.

If
\[
    p(x,y)=-1
\]
on any region of positive area, then probabilities over that region would be negative.
That is also impossible.

\subsection{Average value over a region}

In single-variable calculus, the average value of \(f(x)\) on \([a,b]\) was
\[
    f_{\text{avg}}
    =
    \frac{1}{b-a}\int_a^b f(x)\,dx.
\]
In two variables, we divide by area instead of length.

\begin{definition}[Average value over a region]
If \(f\) is continuous on a region \(D\) with positive area, then the average value of
\(f\) over \(D\) is
\[
\boxed{
    f_{\text{avg}}
    =
    \frac{1}{\operatorname{Area}(D)}
    \iint_D f(x,y)\,dA.
}
\]
\end{definition}

\begin{example}[Average temperature in a circular pond]
A circular pond has radius \(4\) meters.  Its surface temperature depends only on
distance from the center:
\[
    T(r)=20-r.
\]
The center is warmest, and the edge is coolest.

The area of the pond is
\[
    16\pi.
\]
The total temperature integral is
\[
\begin{aligned}
    \iint_D T\,dA
    &=
    \int_0^{2\pi}\int_0^4 (20-r)r\,dr\,d\theta\\
    &=
    2\pi\int_0^4(20r-r^2)\,dr\\
    &=
    2\pi
    \left[
        10r^2-\frac13r^3
    \right]_0^4\\
    &=
    2\pi
    \left(
        160-\frac{64}{3}
    \right)\\
    &=
    \frac{832\pi}{3}.
\end{aligned}
\]
Therefore
\[
\begin{aligned}
    T_{\text{avg}}
    &=
    \frac{1}{16\pi}\cdot\frac{832\pi}{3}\\
    &=
    \frac{52}{3}.
\end{aligned}
\]
So the average temperature is
\[
\boxed{
    \frac{52}{3}^\circ\text{C}\approx 17.33^\circ\text{C}.
}
\]
\end{example}

The same formula can average rainfall, elevation, pollution concentration, light intensity,
or any other scalar field over a region.

\subsection{One integral, many meanings}

The expression
\[
    \iint_D f(x,y)\,dA
\]
has one mathematical structure: add local contributions over area.

Its meaning depends on what \(f\) measures.

\[
\begin{array}{c|c|c}
\text{Quantity} & f(x,y) & \iint_D f(x,y)\,dA \\ \hline
\text{height} & \text{height above a floor} & \text{volume}\\[1mm]
\text{density} & \text{mass per area} & \text{mass}\\[1mm]
\text{charge density} & \text{charge per area} & \text{total charge}\\[1mm]
\text{temperature} & \text{temperature} & \text{area-total used for average}\\[1mm]
\text{probability density} & \text{probability per area} & \text{probability}\\[1mm]
\text{rotational weight} & \text{squared distance}\times\text{density} & \text{moment of inertia}
\end{array}
\]

This is one of the strengths of integration.  Once we know how to add over a region,
many physical and geometric quantities become the same kind of calculation.

\subsection{A compact summary of the main formulas}

Let \(D\) be a region in the plane.

If \(\sigma(x,y)\) is density, then
\[
\boxed{
    m=\iint_D \sigma\,dA.
}
\]

The moments are
\[
\boxed{
    M_y=\iint_D x\sigma\,dA,
    \qquad
    M_x=\iint_D y\sigma\,dA.
}
\]

The center of mass is
\[
\boxed{
    \bar x=\frac{M_y}{m},
    \qquad
    \bar y=\frac{M_x}{m}.
}
\]

The moments of inertia are
\[
\boxed{
    I_x=\iint_D y^2\sigma\,dA,
    \qquad
    I_y=\iint_D x^2\sigma\,dA,
}
\]
and
\[
\boxed{
    I_0=\iint_D (x^2+y^2)\sigma\,dA.
}
\]

If \(p(x,y)\) is a probability density, then
\[
\boxed{
    P((X,Y)\in A)=\iint_A p(x,y)\,dA.
}
\]

If \(f(x,y)\) is a scalar field, then its average value is
\[
\boxed{
    f_{\text{avg}}
    =
    \frac{1}{\operatorname{Area}(D)}
    \iint_D f\,dA.
}
\]

Each formula is a weighted sum.  The weight tells the integral what to measure.

\subsection{What comes next}

So far all our regions have lived in the plane.  We have added over area:
\[
    dA.
\]
But real objects fill space.  A loaf of bread, a storm cloud, a storage tank, and a solid
metal part occupy three-dimensional regions.

The next layer of the story is to add over volume:
\[
    dV.
\]
That leads to triple integrals.

Before that, there is one warning worth understanding.  Some regions are unbounded,
and some functions blow up near a point.  Then the ordinary double integral becomes an
improper double integral.  The same idea of adding tiny pieces remains, but convergence
becomes the new question.
\section{Improper double integrals}
\label{sec:10.7}

Section \ref{sec:10.6} used double integrals to add mass, moments, probability,
rainfall, temperature, and elevation over finite regions.  The pieces were tiny, but the
region itself was bounded, and the function stayed finite.

Some real quantities refuse to be that polite.  Smoke can spread across a whole plain.
A gravitational field can become very large near a point.  A probability density may
live on the entire plane.  To handle these cases, we keep the same idea of adding tiny
pieces, but the final answer now comes from a limit.

\subsection{Unbounded regions}

Imagine a thin layer of dust settling on a flat desert after a volcanic eruption.  The dust
is thickest near the source and thinner farther away.  A simple radial model is
\[
    \rho(x,y)=e^{-\sqrt{x^2+y^2}},
\]
measured in grams per square meter.

The dust has no sharp edge in the model.  It covers the whole plane.  Still, the total
mass might be finite because the density decreases fast.

Use polar coordinates:
\[
    r=\sqrt{x^2+y^2}.
\]
Then
\[
    \rho=e^{-r},
    \qquad
    dA=r\,dr\,d\theta.
\]
The mass inside a disk of radius \(A\) is
\[
\begin{aligned}
M(A)
&=
\int_0^{2\pi}\int_0^A e^{-r}r\,dr\,d\theta\\
&=
2\pi\int_0^A re^{-r}\,dr.
\end{aligned}
\]
Integrate by parts:
\[
    \int re^{-r}\,dr=-re^{-r}-e^{-r}.
\]
So
\[
\begin{aligned}
M(A)
&=
2\pi\left[-re^{-r}-e^{-r}\right]_0^A\\
&=
2\pi\left(1-(A+1)e^{-A}\right).
\end{aligned}
\]
As
\[
    A\to\infty,
\]
we have
\[
    (A+1)e^{-A}\to0.
\]
Therefore
\[
\boxed{
    M=\lim_{A\to\infty}M(A)=2\pi.
}
\]

The dust covers an infinite plane, but its total mass is finite.

\begin{definition}[Improper double integral over an unbounded region]
Let \(R\) be an unbounded region in the plane.  For \(A>0\), let
\[
    R_A=R\cap\{(x,y):x^2+y^2\le A^2\}.
\]
If the limit
\[
\boxed{
    \lim_{A\to\infty}\iint_{R_A} f\,dA
}
\]
exists, then we call it the \emph{improper double integral} of \(f\) over \(R\), and we
write
\[
\boxed{
    \iint_R f\,dA.
}
\]
\end{definition}

For nonnegative functions, this definition matches the physical picture: measure more
and more of the region, then see whether the accumulated total settles to a finite
number.

\begin{figure}[h]
\centering
\begin{tikzpicture}[scale=0.85, >=stealth]
    \begin{scope}
        \draw[->] (-3.4,0) -- (3.7,0) node[right] {\(x\)};
        \draw[->] (0,-3.4) -- (0,3.7) node[above] {\(y\)};

        \foreach \r in {1,2,3} {
            \draw[gray!45] (0,0) circle (\r);
        }

        \draw[->, thick] (2.2,2.2) -- (3.0,3.0);
        \node at (0,-4.0) {expand the disk: \(A\to\infty\)};
    \end{scope}

    \begin{scope}[shift={(8,0)}]
        \draw[->] (-3.2,0) -- (3.5,0) node[right] {\(x\)};
        \draw[->] (0,-3.2) -- (0,3.5) node[above] {\(y\)};

        \draw[fill=green!8, thick] (0,0) circle (2.7);
        \draw[fill=white, thick] (0,0) circle (0.65);

        \draw[->, thick] (0.65,0.65) -- (0.25,0.25);
        \node at (0,-4.0) {shrink the missing disk: \(\varepsilon\to0\)};
    \end{scope}
\end{tikzpicture}
\caption{Improper integrals use limits.  For an unbounded region, the outer boundary
moves outward.  For a singular point, a small excluded disk shrinks inward.}
\end{figure}

The infinite size of the region does not automatically make the integral diverge.  What
matters is the balance between area growth and decay of the function.

Here is the simplest failure.

\begin{example}[Infinite area gives infinite integral]
Let
\[
    f(x,y)=1
\]
on the whole plane.  The integral over the disk of radius \(A\) is
\[
    \iint_{x^2+y^2\le A^2}1\,dA=\pi A^2.
\]
As
\[
    A\to\infty,
\]
this grows without bound.  So
\[
\boxed{
    \iint_{\mathbb R^2}1\,dA
}
\]
does not converge to a finite number.

The plane has infinite area, so adding \(1\) over the whole plane gives an infinite total.
\end{example}

A function can tend to \(0\) and still have infinite total.

\begin{example}[Decay that is too slow]
Let
\[
    f(x,y)=\frac{1}{1+x^2+y^2}
\]
on the whole plane.  In polar coordinates,
\[
    f=\frac{1}{1+r^2}.
\]
The integral over the disk of radius \(A\) is
\[
\begin{aligned}
\int_0^{2\pi}\int_0^A \frac{1}{1+r^2}r\,dr\,d\theta
&=
2\pi\int_0^A \frac{r}{1+r^2}\,dr\\
&=
\pi\ln(1+A^2).
\end{aligned}
\]
As
\[
    A\to\infty,
\]
the logarithm grows without bound.  Therefore
\[
\boxed{
    \iint_{\mathbb R^2}\frac{1}{1+x^2+y^2}\,dA
}
\]
diverges.
\end{example}

The density
\[
    \frac{1}{1+r^2}
\]
goes to \(0\), but not fast enough.  The rings farther and farther away have larger
circumference, and their combined contribution never settles down.

\subsection{Unbounded functions}

A different problem happens when the region is bounded, but the function becomes
infinite at a point.

Imagine heat concentrated near the center of a thin circular plate.  Let the plate be
the unit disk
\[
    D=\{(x,y):x^2+y^2\le1\}.
\]
Suppose the heat density is modeled by
\[
    h(x,y)=\frac{1}{\sqrt{x^2+y^2}}.
\]
This function is unbounded at the origin.  But the origin is only one point, so the
total heat may still be finite.

In polar coordinates,
\[
    h=\frac1r,
    \qquad
    dA=r\,dr\,d\theta.
\]
The product is simple:
\[
    h\,dA=\frac1r\cdot r\,dr\,d\theta=dr\,d\theta.
\]
To avoid the singularity at first, remove a tiny disk of radius \(\varepsilon\).  Then
\[
\begin{aligned}
\iint_{\varepsilon\le r\le1}\frac1r\,dA
&=
\int_0^{2\pi}\int_\varepsilon^1 1\,dr\,d\theta\\
&=
2\pi(1-\varepsilon).
\end{aligned}
\]
Let
\[
    \varepsilon\to0^+.
\]
Then
\[
\boxed{
    \iint_D\frac{1}{\sqrt{x^2+y^2}}\,dA=2\pi.
}
\]
The function blows up at the center, but the total is finite.

\begin{definition}[Improper double integral near a singular point]
Let \(R\) be a bounded region, and suppose \(f\) is unbounded at a point
\[
    \mathbf p\in R.
\]
Let
\[
    R_\varepsilon=R\setminus B_\varepsilon(\mathbf p),
\]
where \(B_\varepsilon(\mathbf p)\) is the disk of radius \(\varepsilon\) centered at
\(\mathbf p\).  If the limit
\[
\boxed{
    \lim_{\varepsilon\to0^+}\iint_{R_\varepsilon} f\,dA
}
\]
exists, then it is called the improper double integral of \(f\) over \(R\).
\end{definition}

The power of \(r\) decides many singularities.  Compare the previous example with
\[
    k(x,y)=\frac{1}{x^2+y^2}.
\]
In polar coordinates,
\[
    k=\frac1{r^2}.
\]
On the punctured disk,
\[
\begin{aligned}
\iint_{\varepsilon\le r\le1}\frac1{r^2}\,dA
&=
\int_0^{2\pi}\int_\varepsilon^1 \frac1{r^2}r\,dr\,d\theta\\
&=
2\pi\int_\varepsilon^1\frac1r\,dr\\
&=
2\pi\ln\frac1{\varepsilon}.
\end{aligned}
\]
As
\[
    \varepsilon\to0^+,
\]
this grows without bound.  So
\[
\boxed{
    \iint_D\frac{1}{x^2+y^2}\,dA
}
\]
diverges.

The difference between
\[
    \frac1r
    \qquad\text{and}\qquad
    \frac1{r^2}
\]
is not small.  The area element contributes a factor of \(r\), and that factor saves
\[
    \frac1r
\]
but not
\[
    \frac1{r^2}.
\]

\subsection{Convergence and comparison}

Most improper double integrals are not computed exactly.  We often only need to know
whether the total is finite.

Suppose a chemical concentration on the whole plane satisfies
\[
    0\le c(x,y)\le \frac{3}{(1+x^2+y^2)^2}.
\]
The upper bound is easier to integrate.  In polar coordinates,
\[
\begin{aligned}
\iint_{\mathbb R^2}\frac{3}{(1+x^2+y^2)^2}\,dA
&=
\int_0^{2\pi}\int_0^\infty \frac{3r}{(1+r^2)^2}\,dr\,d\theta.
\end{aligned}
\]
Use
\[
    u=1+r^2,
    \qquad
    du=2r\,dr.
\]
Then
\[
\begin{aligned}
\int_0^\infty \frac{3r}{(1+r^2)^2}\,dr
&=
\frac32\int_1^\infty u^{-2}\,du\\
&=
\frac32.
\end{aligned}
\]
Therefore
\[
    \iint_{\mathbb R^2}\frac{3}{(1+x^2+y^2)^2}\,dA
    =
    3\pi.
\]
Since \(c\) is squeezed below a function with finite integral, the total amount of
chemical is also finite.

\begin{theorem}[Comparison test for nonnegative double integrals]
Let
\[
    0\le f(x,y)\le g(x,y)
\]
on a region \(R\), except possibly at boundary points or singular points handled by
improper limits.

\begin{enumerate}
\item If
\[
    \iint_R g\,dA
\]
converges, then
\[
    \iint_R f\,dA
\]
also converges.

\item If
\[
    0\le g(x,y)\le f(x,y)
\]
and
\[
    \iint_R g\,dA
\]
diverges, then
\[
    \iint_R f\,dA
\]
also diverges.
\end{enumerate}
\end{theorem}

\begin{proof}[Idea of proof]
For nonnegative functions, the accumulated integral over larger and larger pieces can
only increase.  If it is always bounded above by a finite number, it must settle to a
finite limit.  If it is always at least as large as a quantity that grows without bound,
then it must also grow without bound.
\end{proof}

The nonnegative condition is there because signed functions can cancel.  Cancellation
may be real, but it can also hide size.

For example, over the exterior region \(r\ge1\), the radial function
\[
    f(r,\theta)=\frac{\cos r}{r}
\]
has annular integral
\[
    \int_0^{2\pi}\int_1^A \frac{\cos r}{r}\,r\,dr\,d\theta
    =
    2\pi\int_1^A \cos r\,dr
    =
    2\pi(\sin A-\sin 1).
\]
This has no limit as \(A\to\infty\).  It keeps oscillating.  Positive comparison tests
are designed for totals that do not cancel.

A useful rule of thumb comes from powers of \(r\).

Near the origin, on a disk,
\[
    \frac1{r^p}
\]
has integral like
\[
    \int_0^1 r^{1-p}\,dr.
\]
This converges exactly when
\[
    p<2.
\]

Far from the origin, on the outside of a large disk,
\[
    \frac1{r^p}
\]
has integral like
\[
    \int_1^\infty r^{1-p}\,dr.
\]
This converges exactly when
\[
    p>2.
\]

So in two dimensions:
\[
\begin{array}{c|c}
\text{place} & \dfrac1{r^p}\text{ is integrable when} \\ \hline
\text{near }r=0 & p<2\\[2mm]
\text{near }r=\infty & p>2
\end{array}
\]

The number \(2\) appears because area in polar coordinates carries the extra factor
\[
    r.
\]
In three dimensions, that number will change.

\subsection{The Gaussian integral}

One of the most famous single-variable integrals has no elementary antiderivative:
\[
    \int e^{-x^2}\,dx.
\]
But its improper integral over the whole real line has an exact value.  The surprise is
that a one-dimensional integral becomes easier after we square it and turn it into a
double integral.

Let
\[
    I=\int_{-\infty}^{\infty}e^{-x^2}\,dx.
\]
Since the integrand is positive, we may write
\[
\begin{aligned}
I^2
&=
\left(\int_{-\infty}^{\infty}e^{-x^2}\,dx\right)
\left(\int_{-\infty}^{\infty}e^{-y^2}\,dy\right)\\
&=
\iint_{\mathbb R^2} e^{-x^2}e^{-y^2}\,dA\\
&=
\iint_{\mathbb R^2} e^{-(x^2+y^2)}\,dA.
\end{aligned}
\]
Now use polar coordinates:
\[
    x^2+y^2=r^2,
    \qquad
    dA=r\,dr\,d\theta.
\]
Then
\[
\begin{aligned}
I^2
&=
\int_0^{2\pi}\int_0^\infty e^{-r^2}r\,dr\,d\theta\\
&=
2\pi\int_0^\infty e^{-r^2}r\,dr.
\end{aligned}
\]
Use
\[
    u=r^2,
    \qquad
    du=2r\,dr.
\]
Then
\[
\begin{aligned}
2\pi\int_0^\infty e^{-r^2}r\,dr
&=
2\pi\cdot\frac12\int_0^\infty e^{-u}\,du\\
&=
\pi.
\end{aligned}
\]
Therefore
\[
    I^2=\pi.
\]
Since \(I>0\),
\[
\boxed{
    \int_{-\infty}^{\infty}e^{-x^2}\,dx=\sqrt{\pi}.
}
\]

A single-variable antiderivative was missing, but the double integral found the value.
The plane gave the problem a symmetry that the line could not show.

That is a fitting ending for double integrals.  They began as volume under a surface,
then became a way to add density, probability, and mass across a region.  Even a
one-dimensional integral can become clearer after we let it spread into the plane.
Before moving to space, it is worth gathering the whole chapter into one map.

\section{Chapter review and applications}
\label{sec:10.8}

Chapter \ref{ch:10} began with the same idea that built single-variable integration:
cut, multiply, add, and pass to a limit.  The interval became a region.  The tiny length
\(dx\) became a tiny area \(dA\), or, when orientation matters,
\[
    dx\wedge dy.
\]

The hardest part of most double integral problems is not the integration.  It is deciding
what the small pieces look like.

\subsection{The chapter in one decision path}

A double integral begins with a question:
\[
    \text{What is being added, and over what region?}
\]

If the region is a rectangle, iterated integrals usually go straight through.  If the
region has curved edges, draw it first.  If the region is circular, polar coordinates may
turn the picture into simple bounds.  If the region is infinite or the function blows up,
use an improper limit.

\begin{figure}[h]
\centering
\begin{tikzpicture}[scale=0.88, >=stealth]
    \node[draw, rounded corners, align=center, minimum width=4.2cm, minimum height=0.9cm]
        (start) at (0,0) {Double integral\\ \(\displaystyle \iint_R f\,dA\)};

    \node[draw, rounded corners, align=center, minimum width=4.2cm, minimum height=0.9cm]
        (rect) at (-5,-2) {Rectangle\\ use Fubini};

    \node[draw, rounded corners, align=center, minimum width=4.2cm, minimum height=0.9cm]
        (gen) at (0,-2) {Curved region\\ draw bounds};

    \node[draw, rounded corners, align=center, minimum width=4.2cm, minimum height=0.9cm]
        (polar) at (5,-2) {Circular geometry\\ try polar};

    \node[draw, rounded corners, align=center, minimum width=4.2cm, minimum height=0.9cm]
        (app) at (-2.6,-4.4) {Density, mass,\\ average, probability};

    \node[draw, rounded corners, align=center, minimum width=4.2cm, minimum height=0.9cm]
        (improper) at (2.6,-4.4) {Infinite or singular\\ take a limit};

    \draw[->, thick] (start) -- (rect);
    \draw[->, thick] (start) -- (gen);
    \draw[->, thick] (start) -- (polar);
    \draw[->, thick] (gen) -- (app);
    \draw[->, thick] (polar) -- (app);
    \draw[->, thick] (rect) -- (app);
    \draw[->, thick] (app) -- (improper);
\end{tikzpicture}
\caption{Most double integral problems are solved by matching the region to the right
description.}
\end{figure}

Here are the main formulas from the chapter.

\[
\begin{array}{c|c}
\text{Quantity} & \text{Double integral} \\ \hline
\text{area of }R
&
\displaystyle \operatorname{Area}(R)=\iint_R 1\,dA
\\[3mm]
\text{volume under }z=f(x,y)
&
\displaystyle V=\iint_R f(x,y)\,dA
\\[3mm]
\text{average value}
&
\displaystyle f_{\mathrm{avg}}
=
\frac{1}{\operatorname{Area}(R)}\iint_R f\,dA
\\[4mm]
\text{mass of a lamina}
&
\displaystyle M=\iint_R \rho\,dA
\\[3mm]
\text{moments}
&
\displaystyle M_y=\iint_R x\rho\,dA,
\qquad
M_x=\iint_R y\rho\,dA
\\[4mm]
\text{center of mass}
&
\displaystyle
\bar x=\frac{M_y}{M},
\qquad
\bar y=\frac{M_x}{M}
\\[4mm]
\text{moment of inertia about the origin}
&
\displaystyle
I_O=\iint_R (x^2+y^2)\rho\,dA
\\[4mm]
\text{probability over }R
&
\displaystyle
P((X,Y)\in R)=\iint_R p(x,y)\,dA
\end{array}
\]

For polar coordinates,
\[
\boxed{
    x=r\cos\theta,
    \qquad
    y=r\sin\theta,
    \qquad
    dA=r\,dr\,d\theta.
}
\]
The extra factor
\[
    r
\]
is the area-scaling factor for a small polar rectangle.

For improper double integrals, the same symbols are used, but the meaning includes a
limit:
\[
\boxed{
    \iint_R f\,dA
    =
    \lim_{A\to\infty}\iint_{R_A}f\,dA
}
\]
for an unbounded region, or
\[
\boxed{
    \iint_R f\,dA
    =
    \lim_{\varepsilon\to0^+}\iint_{R_\varepsilon}f\,dA
}
\]
when a singular point has been removed and then allowed back in by a limit.

\subsection{Worked review example: drawing the region before integrating}

Let \(R\) be the region bounded by
\[
    y=x^2
    \qquad\text{and}\qquad
    y=2x.
\]
Compute
\[
    \iint_R y\,dA.
\]

First find the intersection points:
\[
    x^2=2x
    \quad\Longrightarrow\quad
    x(x-2)=0.
\]
So
\[
    x=0
    \qquad\text{or}\qquad
    x=2.
\]
On
\[
    0\le x\le2,
\]
the line \(y=2x\) lies above the parabola \(y=x^2\).  Therefore
\[
    R=\{(x,y):0\le x\le2,\ x^2\le y\le 2x\}.
\]
Set up the integral:
\[
\begin{aligned}
\iint_R y\,dA
&=
\int_0^2\int_{x^2}^{2x} y\,dy\,dx\\
&=
\int_0^2
\left[\frac12y^2\right]_{x^2}^{2x}\,dx\\
&=
\frac12\int_0^2(4x^2-x^4)\,dx\\
&=
\frac12\left[\frac43x^3-\frac15x^5\right]_0^2\\
&=
\frac12\left(\frac{32}{3}-\frac{32}{5}\right)\\
&=
\boxed{\frac{32}{15}}.
\end{aligned}
\]

The same region can be described the other way.  Since
\[
    y=2x
    \quad\Longrightarrow\quad
    x=\frac y2,
\]
and
\[
    y=x^2
    \quad\Longrightarrow\quad
    x=\sqrt y
\]
on the right-hand branch, the reversed description is
\[
    R=\left\{(x,y):0\le y\le4,\ \frac y2\le x\le \sqrt y\right\}.
\]
So the same integral is
\[
\boxed{
    \iint_R y\,dA
    =
    \int_0^4\int_{y/2}^{\sqrt y}y\,dx\,dy.
}
\]
This is the same quantity with a different slicing direction.

\subsection{Worked review example: a circular sprinkler pattern}

A sprinkler waters a circular garden of radius \(2\) meters.  The water density at a
point depends only on distance from the center:
\[
    w(r)=10-r,
    \qquad
    0\le r\le2.
\]
Find the total water and the average water density over the garden.

The region is a disk, and the density is radial, so use polar coordinates:
\[
    0\le r\le2,
    \qquad
    0\le\theta\le2\pi.
\]
The total water is
\[
\begin{aligned}
W
&=
\int_0^{2\pi}\int_0^2 (10-r)r\,dr\,d\theta\\
&=
2\pi\int_0^2(10r-r^2)\,dr\\
&=
2\pi\left[5r^2-\frac13r^3\right]_0^2\\
&=
2\pi\left(20-\frac83\right)\\
&=
\boxed{\frac{104\pi}{3}}.
\end{aligned}
\]
The garden area is
\[
    \pi(2)^2=4\pi.
\]
Therefore the average water density is
\[
\boxed{
    w_{\mathrm{avg}}
    =
    \frac{W}{4\pi}
    =
    \frac{26}{3}.
}
\]

The value
\[
    \frac{26}{3}
\]
is less than \(10\), because most of the garden is not near the center.  In polar
coordinates, the factor \(r\) reminds us that larger rings contain more area.

\subsection{Concept check}

\begin{enumerate}
\item What does
\[
    dA
\]
mean in a double integral?

\item Why is
\[
    \iint_R 1\,dA
\]
the area of \(R\)?

\item What does the inner integral mean geometrically in an iterated integral?

\item Why does changing the order of integration require redrawing the region?

\item What is a Type I region?  What is a Type II region?

\item Why does polar area have the factor
\[
    r
\]
in
\[
    dA=r\,dr\,d\theta?
\]

\item When is polar integration likely to simplify a problem?

\item How are mass and density related by a double integral?

\item Why are the moments
\[
    M_y=\iint_R x\rho\,dA,
    \qquad
    M_x=\iint_R y\rho\,dA
\]
named with the opposite coordinate?

\item What does the average value formula say in words?

\item How can a function be unbounded at one point but still have a finite double
integral?

\item Why can a function that tends to \(0\) at infinity still have a divergent integral?

\item What does the comparison test require?
\end{enumerate}

\subsection{Skill problems}

\begin{enumerate}
\item Compute
\[
    \int_0^3\int_0^2 (x+2y)\,dy\,dx.
\]

\item Compute
\[
    \int_0^1\int_0^x (x^2+y)\,dy\,dx.
\]

\item Let \(R\) be the triangle with vertices
\[
    (0,0),\qquad (4,0),\qquad (4,2).
\]
Set up
\[
    \iint_R f(x,y)\,dA
\]
as an iterated integral in both orders.

\item Let \(R\) be bounded by
\[
    y=x^2
    \qquad\text{and}\qquad
    y=4.
\]
Set up and evaluate
\[
    \iint_R x^2\,dA.
\]

\item Reverse the order of integration:
\[
    \int_0^1\int_y^1 f(x,y)\,dx\,dy.
\]

\item Reverse the order of integration:
\[
    \int_0^4\int_{\sqrt y}^{2} f(x,y)\,dx\,dy.
\]

\item Find the area of the region inside the circle
\[
    x^2+y^2=9
\]
and above the line
\[
    y=0
\]
using polar coordinates.

\item Evaluate
\[
    \iint_D (x^2+y^2)\,dA,
\]
where \(D\) is the disk
\[
    x^2+y^2\le 4.
\]

\item Evaluate
\[
    \iint_A \frac{1}{1+x^2+y^2}\,dA,
\]
where \(A\) is the annulus
\[
    1\le x^2+y^2\le 9.
\]

\item A lamina occupies the rectangle
\[
    0\le x\le2,
    \qquad
    0\le y\le3,
\]
with density
\[
    \rho(x,y)=1+x.
\]
Find its mass.

\item For the lamina in Problem 10, find
\[
    M_x,\qquad M_y,
\]
and its center of mass.

\item A probability density on the square
\[
    0\le x\le1,\qquad 0\le y\le1
\]
is
\[
    p(x,y)=C(x+y).
\]
Find \(C\), and then find the probability that
\[
    x+y\le1.
\]

\item Find the average value of
\[
    f(x,y)=x+y
\]
over the triangle with vertices
\[
    (0,0),\qquad (1,0),\qquad (0,1).
\]

\item Decide whether the improper integral converges:
\[
    \iint_{\mathbb R^2}\frac{1}{(1+x^2+y^2)^2}\,dA.
\]

\item Decide whether the improper integral converges:
\[
    \iint_{\mathbb R^2}\frac{1}{1+x^2+y^2}\,dA.
\]

\item Decide whether the improper integral over the unit disk converges:
\[
    \iint_{x^2+y^2\le1}\frac{1}{(x^2+y^2)^{3/4}}\,dA.
\]

\item Decide whether the improper integral over the unit disk converges:
\[
    \iint_{x^2+y^2\le1}\frac{1}{x^2+y^2}\,dA.
\]

\item Use the Gaussian integral result to evaluate
\[
    \int_{-\infty}^{\infty} e^{-4x^2}\,dx.
\]
\end{enumerate}

\subsection{Project: estimating rainfall from gridded data}

A storm passes over a rectangular watershed.  The watershed is
\[
    0\le x\le8,
    \qquad
    0\le y\le6,
\]
where \(x\) and \(y\) are measured in kilometers.

The region is divided into rectangles of size
\[
    2\text{ km}\times 2\text{ km}.
\]
Rainfall is measured at the center of each rectangle, in centimeters.

\[
\begin{array}{c|cccc}
 & x=1 & x=3 & x=5 & x=7 \\ \hline
y=5 & 3.2 & 4.1 & 4.6 & 3.8\\
y=3 & 2.4 & 3.5 & 5.0 & 4.2\\
y=1 & 1.1 & 2.0 & 3.1 & 2.7
\end{array}
\]

Each grid cell has area
\[
    4\text{ km}^2.
\]

\subsubsection*{Part 1. Estimate the total rainfall integral}

Use a midpoint sum:
\[
    \iint_R r(x,y)\,dA
    \approx
    4\sum r(x_i^*,y_j^*).
\]
For the data above,
\[
\begin{aligned}
\sum r(x_i^*,y_j^*)
&=
3.2+4.1+4.6+3.8\\
&\qquad
+2.4+3.5+5.0+4.2\\
&\qquad
+1.1+2.0+3.1+2.7\\
&=
39.7.
\end{aligned}
\]
So
\[
\boxed{
    \iint_R r(x,y)\,dA
    \approx
    4(39.7)=158.8
    \text{ cm}\cdot\text{km}^2.
}
\]

\subsubsection*{Part 2. Convert to water volume}

One centimeter of rain over one square kilometer equals
\[
    10{,}000\text{ m}^3
\]
of water.  Therefore the estimated water volume is
\[
\boxed{
    158.8(10{,}000)=1{,}588{,}000\text{ m}^3.
}
\]

\subsubsection*{Part 3. Estimate average rainfall}

The watershed area is
\[
    8\cdot6=48\text{ km}^2.
\]
The average rainfall is
\[
\boxed{
    r_{\mathrm{avg}}
    \approx
    \frac{158.8}{48}
    \approx
    3.31\text{ cm}.
}
\]

\subsubsection*{Part 4. What to hand in}

Write a short report with the following pieces.

\begin{enumerate}
\item A sketch of the watershed grid.

\item The midpoint-sum estimate for
\[
    \iint_R r(x,y)\,dA.
\]

\item The conversion from
\[
    \text{cm}\cdot\text{km}^2
\]
to cubic meters.

\item The estimated average rainfall.

\item A contour-style sketch showing where the rainfall appears largest.

\item A paragraph explaining one possible source of error in the estimate.  For example,
what if the heaviest rain fell between two sensors?
\end{enumerate}

A double integral sees a thin region: a floor, a plate, a map, a lamina.  But many
objects are not thin.  Air fills a room.  Charge fills a solid ball.  Density can vary from
one shelf of a tank to another.  The same adding idea is waiting there too, but the tiny
pieces are no longer small rectangles.  They are small boxes, and the bounds must
describe a three-dimensional body.  One more coordinate turns the double integral into
a triple integral.